% Auto-split: keep each TeX source comfortably below ~600 lines.
%
% This subsubsection supplements the multiplicity-composition thermodynamics with:
% - exact finite-volume conditional i.i.d. representations,
% - a soft-core chain (transfer-matrix / 传递矩阵) realization for integer replica index q,
% - explicit S_q symmetry reduction and an auditable φ-power eigen-spectrum,
% - Palm size-bias limits and exponential mixing constants.
%

\subsubsection{纤维权有序分解的精确更新与软核链等价：条件独立、谱隙与 \texorpdfstring{$\varphi$}{phi}-幂级谱}\label{subsec:pom-multiplicity-composition-exact-update-softcore}
沿用定义 \ref{def:pom-multiplicity-composition-partition}--\ref{def:pom-multiplicity-real-q-pressure} 的记号：
\(\mathcal{C}_L\) 为总长度 \(L\) 的全体有序分解，
\(
w_q(\boldsymbol{\ell})=\prod_i F_{\ell_i+2}^{\,q}
\),
\(
Z_L^{(q)}=\sum_{\boldsymbol{\ell}\in\mathcal{C}_L}w_q(\boldsymbol{\ell})
\),
\(
P_L^{(q)}(\boldsymbol{\ell})=w_q(\boldsymbol{\ell})/Z_L^{(q)}
\)
（当 \(q=1\) 时 \(P_L^{(1)}=P_L=\mathbb{P}_L\)，并且 \(\lambda(1)=\lambda_1=\frac{3+\sqrt{17}}{2}\)、\(\rho_\star=\lambda(1)^{-1}=\frac{\sqrt{17}-3}{4}\)）。
本节给出两类互补的严格等价表示：其一为每个有限 \(L\) 的精确条件独立更新链；
其二为整数 \(q\) 下的 \(\{0,1\}^q\) 软核链传递矩阵表示及其对称降维，
并由此导出可审计的谱隙常数、Palm 尺度偏置与指数混合。

\begin{theorem}[纤维权分解的精确条件独立更新表示：有限体积严格等价]\label{thm:pom-multiplicity-composition-exact-conditional-iid}
固定实数 \(q\ge 0\)。
令 \(\rho(q)\in(0,1)\) 为定义 \ref{def:pom-multiplicity-real-q-pressure} 中满足
\(
\sum_{k\ge 1}F_{k+2}^{\,q}\rho(q)^k=1
\)
的唯一解，并定义更新分布
\[
p_q(k):=F_{k+2}^{\,q}\rho(q)^k\qquad(k\ge 1),
\qquad
\sum_{k\ge 1}p_q(k)=1.
\]
令 \((\ell_1,\ell_2,\dots)\) 为独立同分布增量、\(\ell_i\sim p_q\)，并记部分和 \(S_n=\sum_{i=1}^n\ell_i\) 以及首次命中时刻
\[
\tau_L:=\inf\{n\ge 1:\ S_n=L\}.
\]
则对每个 \(L\ge 1\) 与任意 \(\boldsymbol{\ell}=(\ell_1,\dots,\ell_r)\in\mathcal{C}_L\) 有严格恒等式
\begin{equation}\label{eq:pom-multiplicity-composition-exact-conditional-iid}
\boxed{\;
\mathbb{P}\Bigl((\ell_1,\dots,\ell_{\tau_L})=\boldsymbol{\ell}\ \Big|\ \tau_L<\infty\Bigr)
=P_L^{(q)}(\boldsymbol{\ell})
=\frac{w_q(\boldsymbol{\ell})}{Z_L^{(q)}}\; }.
\end{equation}
并且命中概率具有完全闭式
\begin{equation}\label{eq:pom-multiplicity-composition-hit-prob}
\boxed{\;
\mathbb{P}(\tau_L<\infty)=\rho(q)^{L}\,Z_L^{(q)}\; }.
\end{equation}
因此，\(P_L^{(q)}\) 不仅在体极限上呈“局部 i.i.d.”，而是对每个有限 \(L\) 精确等价于一条 i.i.d.\ 更新链在事件 \(\{S_{\tau_L}=L\}\) 上的条件分布。
\leanverified{paper\_pom\_multiplicity\_composition\_exact\_conditional\_iid}
\end{theorem}
\begin{proof}
对任意 \(\boldsymbol{\ell}=(\ell_1,\dots,\ell_r)\in\mathcal{C}_L\)，由独立性与 \(p_q(k)=F_{k+2}^{\,q}\rho(q)^k\) 有
\[
\mathbb{P}\bigl((\ell_1,\dots,\ell_r)=\boldsymbol{\ell}\bigr)
=\prod_{i=1}^r F_{\ell_i+2}^{\,q}\rho(q)^{\ell_i}
=w_q(\boldsymbol{\ell})\,\rho(q)^L.
\]
由于增量严格为正，事件 \(\{\tau_L<\infty\}\) 等价于“某个前缀和恰为 \(L\)”，并由全体有序分解分割为不交并：
\(
\{\tau_L<\infty\}=\bigsqcup_{\boldsymbol{\ell}\in\mathcal{C}_L}\{(\ell_1,\dots,\ell_{\tau_L})=\boldsymbol{\ell}\}
\)。
对上述概率求和即得 \eqref{eq:pom-multiplicity-composition-hit-prob}。
再用条件概率公式得到
\[
\mathbb{P}\Bigl((\ell_1,\dots,\ell_{\tau_L})=\boldsymbol{\ell}\ \Big|\ \tau_L<\infty\Bigr)
=\frac{w_q(\boldsymbol{\ell})\,\rho(q)^L}{\rho(q)^L Z_L^{(q)}}
=\frac{w_q(\boldsymbol{\ell})}{Z_L^{(q)}}
=P_L^{(q)}(\boldsymbol{\ell}),
\]
即为 \eqref{eq:pom-multiplicity-composition-exact-conditional-iid}。证毕。
\end{proof}

\begin{proposition}[首部件长度的有限窗精确公式与指数收敛速率]\label{prop:pom-multiplicity-composition-first-part-exact-and-rate}
\leanverified{paper\_pom\_multiplicity\_composition\_first\_part\_exact\_and\_rate}
固定 \(q\ge 0\)。对任意 \(L\ge 1\) 与 \(1\le k\le L\)，在 \(P_L^{(q)}\) 下首部件长度满足严格恒等式
\begin{equation}\label{eq:pom-multiplicity-composition-first-part-exact}
\boxed{\;
P_L^{(q)}(\ell_1=k)=\frac{F_{k+2}^{\,q}\,Z_{L-k}^{(q)}}{Z_L^{(q)}}\; }.
\end{equation}
此外，当 \(q=1\) 时存在常数 \(C_k<\infty\) 使对所有 \(L\ge 1\) 成立指数界
\begin{equation}\label{eq:pom-multiplicity-composition-first-part-exp-rate}
\boxed{\;
\Bigl|P_L(\ell_1=k)-F_{k+2}\rho_\star^{k}\Bigr|
\le C_k\Bigl(\frac{2}{\lambda(1)^2}\Bigr)^{L}\; } ,
\qquad
\rho_\star=\lambda(1)^{-1}.
\end{equation}
其中 \(\frac{2}{\lambda(1)^2}=\frac{4}{13+3\sqrt{17}}\in(0,1)\) 为显式代数谱隙常数。
\end{proposition}
\begin{proof}
恒等式 \eqref{eq:pom-multiplicity-composition-first-part-exact} 由“首部件长度 \(k\) + 尾部为 \(L-k\) 的任意分解”并用权重的串接乘法性直接得到。
\par
对指数界，只需用命题 \ref{prop:pom-multiplicity-composition-partition-rational} 的闭式
\(
Z_L=a\lambda_+^L+b\lambda_-^L
\)
（其中 \(\lambda_+=\lambda(1)\)、\(a=\frac{17+\sqrt{17}}{34}>0\)、\(b=\frac{17-\sqrt{17}}{34}>0\)、\(\lambda_-=\frac{3-\sqrt{17}}{2}\)）
比较比值 \(Z_{L-k}/Z_L\) 与 \(\lambda_+^{-k}=\rho_\star^{k}\)：
令 \(r:=\lambda_-/\lambda_+\)，则 \(|r|=\frac{2}{\lambda(1)^2}\) 且对所有 \(L\ge 1\) 有
\[
\frac{Z_{L-k}}{Z_L}
=\lambda_+^{-k}\cdot\frac{a+br^{L-k}}{a+br^L},
\qquad
\frac{Z_{L-k}}{Z_L}-\lambda_+^{-k}
=\lambda_+^{-k}\cdot\frac{b(r^{L-k}-r^L)}{a+br^L}.
\]
用下界 \(a+br^L\ge a-b|r|\)（对 \(L\ge 1\)）与 \(|r^{L-k}-r^L|\le |r|^{L-k}(1+|r|^k)\) 得
\[
\Bigl|\frac{Z_{L-k}}{Z_L}-\lambda_+^{-k}\Bigr|
\le
\lambda_+^{-k}\cdot \frac{b(1+|r|^k)}{a-b|r|}\,|r|^{L-k}
\le
\lambda_+^{-k}\cdot \frac{b(1+|r|^k)}{a-b|r|}\,|r|^{-k}\,|r|^L.
\]
乘以 \(F_{k+2}\) 并取
\(
C_k:=F_{k+2}\lambda_+^{-k}\frac{b(1+|r|^k)}{a-b|r|}|r|^{-k}
\)
即可得到 \eqref{eq:pom-multiplicity-composition-first-part-exp-rate}。证毕。
\end{proof}

\begin{algorithm}[Doob--更新核的精确序贯生成：残长马尔可夫链与吸收采样]\label{alg:pom-multiplicity-composition-doob-update-sampler}
固定 \(q\ge 0\) 与 \(L\ge 1\)。定义残长过程 \(N_0=L\)，并对任意 \(N\ge 1\) 定义一步核
\begin{equation}\label{eq:pom-multiplicity-composition-doob-kernel}
\boxed{\;
\mathbb{P}(N\to N-k)=\frac{F_{k+2}^{\,q}\,Z_{N-k}^{(q)}}{Z_N^{(q)}},
\qquad 1\le k\le N\; }.
\end{equation}
令过程在 \(N=0\) 吸收，并把每一步的跳长记为 \(k_1,k_2,\dots,k_{\tau}\)（其中 \(\tau\) 为吸收时间）。
则 \((k_1,\dots,k_\tau)\in\mathcal{C}_L\) 且其分布严格等于 \(P_L^{(q)}\)：
\[
(k_1,\dots,k_\tau)\ \sim\ P_L^{(q)}.
\]
\end{algorithm}
\begin{proof}
由 \eqref{eq:pom-multiplicity-composition-doob-kernel}，
对任意 \(\boldsymbol{k}=(k_1,\dots,k_r)\in\mathcal{C}_L\) 有
\[
\mathbb{P}\bigl((k_1,\dots,k_\tau)=\boldsymbol{k}\bigr)
=\prod_{i=1}^{r}\frac{F_{k_i+2}^{\,q}Z_{L-\sum_{j\le i}k_j}^{(q)}}{Z_{L-\sum_{j<i}k_j}^{(q)}}
=\frac{\prod_{i=1}^{r}F_{k_i+2}^{\,q}}{Z_L^{(q)}}
=P_L^{(q)}(\boldsymbol{k}),
\]
其中分式连乘按望远镜消去化简。证毕。
\end{proof}

\begin{proposition}[整阶矩配分函数的 \texorpdfstring{$q$}{q}-复制软核链表示：状态空间 \texorpdfstring{$\{0,1\}^q$}{\{0,1\}^q} 的传递矩阵刻画]\label{prop:pom-multiplicity-composition-replica-softcore-transfer}
\leanverified{paper\_pom\_multiplicity\_composition\_replica\_softcore\_transfer}
取整数 \(q\ge 1\)。
令 \(U_q:=\{0,1\}^q\) 并把 \(u\wedge v\in\{0,1\}^q\) 记为逐坐标与。
定义传递矩阵 \(T^{(q)}\in\RR^{U_q\times U_q}\) 为
\[
T^{(q)}_{u,v}:=2^{-\mathbf{1}_{\{u\wedge v\neq 0\}}}
\qquad(u,v\in U_q).
\]
则对任意 \(L\ge 1\) 成立严格等价
\begin{equation}\label{eq:pom-multiplicity-composition-replica-transfer}
\boxed{\;
Z_L^{(q)}=2^{L-1}\sum_{(u_1,\dots,u_L)\in U_q^{L}}\ 2^{-\sum_{i=1}^{L-1}\mathbf{1}_{\{u_i\wedge u_{i+1}\neq 0\}}}
=2^{L-1}\,\mathbf{1}^{\mathsf T}\bigl(T^{(q)}\bigr)^{L-1}\mathbf{1}\; } ,
\end{equation}
其中 \(\mathbf{1}\) 为 \(\RR^{U_q}\) 上的全 \(1\) 向量。
从而整数一致性增长常数满足
\begin{equation}\label{eq:pom-multiplicity-composition-lambdaq-transfer-radius}
\boxed{\;
\lambda(q)=\lambda_q=2\,\rho\!\bigl(T^{(q)}\bigr)\qquad(q\in\ZZ_{\ge 1}),\;}
\end{equation}
其中 \(\rho(\cdot)\) 表示谱半径。
\end{proposition}
\begin{proof}
把 \(\boldsymbol{\ell}\in\mathcal{C}_L\) 视为在 \(\{1,\dots,L-1\}\) 上的切边集 \(C(\boldsymbol{\ell})\subseteq\{1,\dots,L-1\}\)：
若 \(i\) 为某个部件的右端点则 \(i\in C(\boldsymbol{\ell})\)（于是 \(|C|=R(\boldsymbol{\ell})-1\)）。
\par
另一方面，\(F_{\ell+2}\) 计数长度 \(\ell\) 上“无相邻 \(1\)”的二元串；
因此 \(F_{\ell+2}^{\,q}\) 计数 \(q\) 份独立的无相邻 \(1\) 串，等价于计数向量串 \((u_1,\dots,u_\ell)\in U_q^\ell\) 满足对每个相邻对 \(u_i\wedge u_{i+1}=0\)（逐坐标无重叠）。
由此，给定切边集 \(C\) 后，分量乘法给出
\[
w_q(\boldsymbol{\ell})
=\prod_{\text{分量}}F_{\ell_i+2}^{\,q}
=\#\Bigl\{(u_1,\dots,u_L)\in U_q^L:\ \forall i\notin C,\ u_i\wedge u_{i+1}=0\Bigr\}.
\]
对固定 \((u_1,\dots,u_L)\)，令违约边集合
\(
S(u):=\{i:\ u_i\wedge u_{i+1}\neq 0\}
\)。
则上述条件等价于 \(C\supseteq S(u)\)，并且满足 \(C\supseteq S(u)\) 的切边集个数为 \(2^{(L-1)-|S(u)|}\)。
交换求和次序得到
\[
Z_L^{(q)}
=\sum_{C\subseteq\{1,\dots,L-1\}}\#\{u:\ C\supseteq S(u)\}
=\sum_{u\in U_q^L}2^{(L-1)-|S(u)|}
=2^{L-1}\sum_{u\in U_q^L}2^{-\sum_{i=1}^{L-1}\mathbf{1}_{\{u_i\wedge u_{i+1}\neq 0\}}},
\]
即得 \eqref{eq:pom-multiplicity-composition-replica-transfer}。
最后，\(Z_L^{(q)}=2^{L-1}\mathbf{1}^{\mathsf T}(T^{(q)})^{L-1}\mathbf{1}\) 的指数增长率由 Perron--Frobenius 给出为
\(
\log 2+\log\rho(T^{(q)})=\log(2\rho(T^{(q)}))
\)，从而得到 \eqref{eq:pom-multiplicity-composition-lambdaq-transfer-radius}。证毕。
\end{proof}

\begin{proposition}[\texorpdfstring{$S_q$}{Sq}-对称降维：\(\rho(T^{(q)})\) 为显式 \texorpdfstring{$(q+1)\times(q+1)$}{(q+1)x(q+1)} 矩阵的 Perron 根]\label{prop:pom-multiplicity-composition-sq-symmetry-reduction}
\leanverified{paper\_pom\_multiplicity\_composition\_sq\_symmetry\_reduction}
在命题 \ref{prop:pom-multiplicity-composition-replica-softcore-transfer} 的整数 \(q\ge 1\) 设定下，
令 \(|u|\) 表示 Hamming 权。
定义显式矩阵 \(A^{(q)}\in\RR^{(q+1)\times(q+1)}\) 为
\begin{equation}\label{eq:pom-multiplicity-composition-Aq-def}
\boxed{\;
A^{(q)}_{k,l}:=\frac12\Bigl(\binom{q}{l}+\binom{q-k}{l}\Bigr),
\qquad 0\le k,l\le q.\;}
\end{equation}
则 Perron 根满足严格等式
\begin{equation}\label{eq:pom-multiplicity-composition-sq-reduction}
\boxed{\;
\rho\!\bigl(T^{(q)}\bigr)=\rho\!\bigl(A^{(q)}\bigr),\qquad
\lambda_q=2\,\rho\!\bigl(A^{(q)}\bigr).\;}
\end{equation}
因此对每个整数 \(q\ge 1\)，增长常数 \(\lambda_q\) 为次数至多 \(q+1\) 的代数数（由 \(\det(\rho(A^{(q)})I-A^{(q)})=0\) 给出）。
\end{proposition}
\begin{proof}
矩阵 \(T^{(q)}\) 在坐标置换群 \(S_q\) 下不变，故其 Perron 特征向量可取为 \(S_q\)-不变的类函数：
存在正函数 \(f:U_q\to\RR_+\) 仅依赖于 \(k=|u|\)。
把 \(f(u)\equiv g(|u|)\) 代入 \((T^{(q)}f)(u)=\sum_{v\in U_q}T^{(q)}_{u,v}f(v)\)，并对 \(v\) 按权 \(l=|v|\) 分组：
对固定 \(u\)（\(|u|=k\)）与固定 \(l\)，满足 \(|v|=l\) 且 \(u\wedge v=0\) 的 \(v\) 个数为 \(\binom{q-k}{l}\)，而权为 \(l\) 的总数为 \(\binom{q}{l}\)。
因此
\[
\sum_{\substack{v\in U_q\\ |v|=l}}T^{(q)}_{u,v}
=\binom{q-k}{l}\cdot 1 + \bigl(\binom{q}{l}-\binom{q-k}{l}\bigr)\cdot \frac12
=\frac12\Bigl(\binom{q}{l}+\binom{q-k}{l}\Bigr)
=A^{(q)}_{k,l}.
\]
这表明 \(T^{(q)}\) 在 \(S_q\)-不变子空间（维数 \(q+1\)）上的矩阵表示正是 \(A^{(q)}\)，且 Perron 本征值位于该子空间并等于 \(\rho(A^{(q)})\)。
由于 \(T^{(q)}\) 逐项正，Perron 根唯一且必由该 \(S_q\)-不变 Perron 向量实现，故 \(\rho(T^{(q)})=\rho(A^{(q)})\)。
结合 \(\lambda_q=2\rho(T^{(q)})\) 得 \eqref{eq:pom-multiplicity-composition-sq-reduction}。证毕。
\end{proof}

\begin{proposition}[\texorpdfstring{$q=1$}{q=1} 的完全可解马尔可夫化：平稳分布、转移核与交替相关长度]\label{prop:pom-multiplicity-composition-q1-markov}
\leanverified{paper\_pom\_multiplicity\_composition\_q1\_markov}
在 \(q=1\) 时，令
\[
T:=T^{(1)}=\begin{pmatrix}1&1\\[2pt]1&\tfrac12\end{pmatrix},
\qquad
\alpha:=\rho(T)=\frac{3+\sqrt{17}}{4},
\qquad
\rho_\star:=\lambda(1)^{-1}=\frac{1}{2\alpha}.
\]
考虑一维 Gibbs 权
\(
\pi_L(\sigma)\propto \prod_{i=1}^{L-1}T_{\sigma_i,\sigma_{i+1}}
\)
（\(\sigma\in\{0,1\}^L\)）。
则其体极限存在且唯一，并等价于二态平移不变马尔可夫链 \((\sigma_i)_{i\in\ZZ}\)，其转移核满足
\[
\boxed{\;
\mathbb{P}(\sigma_{i+1}=1\mid\sigma_i=1)=\rho_\star,\qquad
\mathbb{P}(\sigma_{i+1}=1\mid\sigma_i=0)=1-2\rho_\star.\;}
\]
该链具有平稳分布闭式
\[
\boxed{\;
\mathbb{P}(\sigma_i=1)=\frac{17-\sqrt{17}}{34},\qquad
\mathbb{P}(\sigma_i=0)=\frac{17+\sqrt{17}}{34}.\;}
\]
并且二点相关严格指数衰减且符号交替：对任意仅依赖单点的有界函数 \(f,g\)，存在常数 \(C(f,g)\) 使
\[
\mathrm{Cov}\bigl(f(\sigma_0),g(\sigma_n)\bigr)
=C(f,g)\Bigl(-\frac{2}{\lambda(1)^2}\Bigr)^{n}\qquad(n\ge 0).
\]
因此相关长度由 \(\xi^{-1}=\log(\lambda(1)^2/2)\) 给出。
\end{proposition}
\begin{proof}[证明要点]
由对称转移矩阵 \(T\) 的 Perron--Frobenius 理论，体极限 Gibbs 态由 Perron 特征向量 \(r\) 给出，并且可写为马尔可夫链：
\(
P(a\to b)=T_{ab}r_b/(\alpha r_a)
\)。
直接解 \(Tr=\alpha r\) 得 \(r\propto(1,\alpha-1)\)，代回即得转移核闭式；
平稳分布由可逆性 \(\pi(a)\propto r_a^2\) 化简得到。
相关衰减由 \(2\times 2\) 链的第二特征值比给出；
由于 \(2T\) 的特征值为 \(\lambda_\pm=\frac{3\pm\sqrt{17}}{2}\) 且 \(\lambda_+\lambda_-=-2\)，得到衰减因子
\(
\lambda_-/\lambda_+=-2/\lambda(1)^2
\)
并呈交替符号。证毕。
\end{proof}

\begin{corollary}[切边密度的代数闭式：分解密度与“软核重叠”强迫切断的平均]\label{cor:pom-multiplicity-composition-cut-density}
对 \(L\ge 2\)，把 \(\boldsymbol{\ell}\in\mathcal{C}_L\) 等价为切边变量 \(C_i\in\{0,1\}\)（\(1\le i\le L-1\)）：
当且仅当 \(i\) 为某个部件右端点时 \(C_i=1\)。
则在体极限中切边密度满足闭式
\[
\boxed{\;
\lim_{L\to\infty}\frac{1}{L-1}\mathbb{E}_{P_L}\Bigl[\sum_{i=1}^{L-1}C_i\Bigr]
=\frac{17+5\sqrt{17}}{68}\; }.
\]
从而典型部件长度极限为其倒数
\[
\boxed{\;
\lim_{L\to\infty}\frac{L}{\mathbb{E}_{P_L}[R_L]}
=\frac{5\sqrt{17}-17}{2}\; } ,
\]
其中 \(R_L\) 为部件数（定义 \ref{def:pom-multiplicity-induced-composition-measure}）。
\end{corollary}
\begin{proof}
由命题 \ref{prop:pom-multiplicity-composition-replica-softcore-transfer} 在 \(q=1\) 的表示，
\(Z_L\) 亦可视为对偶变量 \((\sigma,C)\) 的一致计数：对每条边 \(i\)，若 \(\sigma_i=\sigma_{i+1}=1\) 则强迫 \(C_i=1\)，否则 \(C_i\) 自由取值。
因此在 Gibbs 极限下，
\[
\mathbb{E}[C_i]=\frac12+\frac12\,\mathbb{P}(\sigma_i=\sigma_{i+1}=1).
\]
由命题 \ref{prop:pom-multiplicity-composition-q1-markov} 的平稳分布与转移核，
\(
\mathbb{P}(\sigma_i=\sigma_{i+1}=1)=\mathbb{P}(\sigma_i=1)\mathbb{P}(\sigma_{i+1}=1\mid\sigma_i=1)
=\frac{17-\sqrt{17}}{34}\cdot \frac{\sqrt{17}-3}{4}
=\frac{5\sqrt{17}-17}{34}
\)。
代入得
\(
\mathbb{E}[C_i]=\frac12+\frac{5\sqrt{17}-17}{68}=\frac{17+5\sqrt{17}}{68}
\)。
又因 \(R_L=1+\sum_{i=1}^{L-1}C_i\)，除以 \(L\) 并取极限得到第二式。证毕。
\end{proof}

\begin{proposition}[Palm 尺度的部件长度分布：均匀位置看到的长度为严格尺寸偏置]\label{prop:pom-multiplicity-composition-palm-size-bias}
在 \(P_L\) 下令 \(V_L\) 为从 \(\{1,\dots,L\}\) 均匀取的随机位置，并记 \(\mathcal{L}_L\) 为包含 \(V_L\) 的部件长度。
令 \(p(k)=F_{k+2}\rho_\star^k\) 为定理 \ref{thm:pom-multiplicity-composition-exact-conditional-iid} 在 \(q=1\) 的更新分布，且记
\(
\mathbb{E}[\ell]=\sum_{k\ge 1}k\,p(k)=\frac{5\sqrt{17}-17}{2}
\)（命题 \ref{prop:pom-multiplicity-composition-part-length-iid-limit}）。
则当 \(L\to\infty\) 时，\(\mathcal{L}_L\) 收敛到尺寸偏置极限
\begin{equation}\label{eq:pom-multiplicity-composition-palm-size-bias-law}
\boxed{\;
\mathbb{P}(\mathcal{L}=k)=\frac{k\,p(k)}{\mathbb{E}[\ell]}
=\frac{k\,F_{k+2}\rho_\star^k}{\mathbb{E}[\ell]}\qquad(k\ge 1).\;}
\end{equation}
并且其均值为闭式
\[
\boxed{\;
\mathbb{E}[\mathcal{L}]
=\frac{\mathbb{E}[\ell^2]}{\mathbb{E}[\ell]}
=\frac{(393-93\sqrt{17})/2}{(5\sqrt{17}-17)/2}
=\frac{393-93\sqrt{17}}{5\sqrt{17}-17}\; } .
\]
\end{proposition}
\begin{proof}[证明要点]
把分解视为更新过程生成的 renewal 分割，并以“均匀位置”诱导 Palm 视角：
在平移不变更新极限中，落在长度为 \(k\) 的部件内的概率与其长度成正比，因此极限分布为尺寸偏置 \(\propto k\,p(k)\)。
均值恒等式 \(\mathbb{E}[\mathcal{L}]=\mathbb{E}[\ell^2]/\mathbb{E}[\ell]\) 为标准 Palm 计算；
其中
\(
\mathbb{E}[\ell^2]=\mathrm{Var}(\ell)+\mathbb{E}[\ell]^2
=(18-4\sqrt{17})+\bigl(\frac{5\sqrt{17}-17}{2}\bigr)^2
=(393-93\sqrt{17})/2
\)
由命题 \ref{prop:pom-multiplicity-composition-part-length-iid-limit} 立即得到。证毕。
\end{proof}

\begin{theorem}[指数混合：宏观分解的远距独立性由谱隙 \texorpdfstring{$2/\lambda(1)^2$}{2/lambda(1)^2} 控制]\label{thm:pom-multiplicity-composition-exponential-mixing}
\leanverified{paper\_pom\_multiplicity\_composition\_exponential\_mixing}
令 \(A,B\) 为仅依赖于切边变量 \((C_i)\) 在两个相距至少 \(n\) 条边的有限窗口上的任意柱事件。
则在 \(L\to\infty\) 的体极限中，存在常数 \(C(A,B)<\infty\) 使
\[
\boxed{\;
\bigl|\mathbb{P}(A\cap B)-\mathbb{P}(A)\mathbb{P}(B)\bigr|
\le C(A,B)\Bigl(\frac{2}{\lambda(1)^2}\Bigr)^{n}\; } .
\]
\end{theorem}
\begin{proof}[证明要点]
由 \(q=1\) 的软核链表示，\((C_i)\) 可由隐藏马尔可夫链 \((\sigma_i)\) 生成：若 \((\sigma_i,\sigma_{i+1})=(1,1)\) 则 \(C_i\equiv 1\)，否则 \(C_i\) 条件独立且 \(\mathbb{P}(C_i=1\mid \sigma_i\sigma_{i+1}=0)=1/2\)。
命题 \ref{prop:pom-multiplicity-composition-q1-markov} 给出 \((\sigma_i)\) 的谱隙与相关衰减因子 \(|-2/\lambda(1)^2|<1\)；
对给定窗口事件 \(A,B\) 先对 \(\sigma\) 条件化并用马尔可夫链的指数混合，再用给定 \(\sigma\) 时 \(C\) 在边上的条件独立性即可把同一衰减率转移到 \(C\) 的柱事件上。证毕。
\end{proof}

\begin{proposition}[\texorpdfstring{$q$}{q}-复制谱的显式非对称本征序列：大量本征值为 \texorpdfstring{$(-\varphi^{-1})$}{-phi^{-1}}-幂的可审计结构]\label{prop:pom-multiplicity-composition-replica-explicit-eigs}
在命题 \ref{prop:pom-multiplicity-composition-replica-softcore-transfer} 的整数 \(q\ge 1\) 设定下，
令
\[
K=\begin{pmatrix}1&1\\[2pt]1&0\end{pmatrix},
\qquad
\text{其特征值为 }\ \varphi=\frac{1+\sqrt5}{2}\ \text{与}\ -\varphi^{-1}.
\]
令 \(D:=K^{\otimes q}\)（按自然张量基识别为 \(U_q\times U_q\) 矩阵），则有严格恒等式
\[
\boxed{\;
T^{(q)}=\frac12\bigl(\mathbf{1}\mathbf{1}^{\mathsf T}+D\bigr),\qquad
D_{u,v}=\mathbf{1}_{\{u\wedge v=0\}}\; } .
\]
因此对每个 \(0\le j\le q\)，数
\begin{equation}\label{eq:pom-multiplicity-composition-replica-explicit-eigs}
\boxed{\;
\frac{(-1)^j\varphi^{\,q-2j}}{2}\;}
\end{equation}
作为 \(T^{(q)}\) 的本征值至少具有重数 \(\binom{q}{j}-1\)（在与 \(\mathbf{1}\) 正交的部分）。
其余至多 \(q+1\) 个本征值全部落在 \(S_q\)-对称子空间，并由命题 \ref{prop:pom-multiplicity-composition-sq-symmetry-reduction} 的显式矩阵 \(A^{(q)}\) 给出。
\end{proposition}
\begin{proof}
对 \(u,v\in\{0,1\}\)，有
\(
K_{u,v}=\mathbf{1}_{\{uv=0\}}
\)。
由张量积的逐坐标乘法性，
\(
D_{u,v}=\prod_{a=1}^q \mathbf{1}_{\{u_av_a=0\}}=\mathbf{1}_{\{u\wedge v=0\}}
\)；
于是
\(
T^{(q)}_{u,v}=1\) 当且仅当 \(u\wedge v=0\)，否则 \(T^{(q)}_{u,v}=1/2\)，即 \(T^{(q)}=\frac12(\mathbf{1}\mathbf{1}^{\mathsf T}+D)\)。
\par
矩阵 \(D=K^{\otimes q}\) 的本征值为张量谱的乘积：
对每个 \(j\)，取 \(j\) 个因子给出本征值 \(-\varphi^{-1}\)、其余 \(q-j\) 个因子取 \(\varphi\)，得到本征值 \((-1)^j\varphi^{q-2j}\)，其代数重数为 \(\binom{q}{j}\)。
另一方面 \(\mathbf{1}\mathbf{1}^{\mathsf T}\) 在 \(\mathbf{1}^\perp\) 上为零，故在 \(\mathbf{1}^\perp\) 上有 \(T^{(q)}|_{\mathbf{1}^\perp}=\frac12 D|_{\mathbf{1}^\perp}\)。
任一 \(\binom{q}{j}\)-维本征空间与 \(\mathbf{1}^\perp\)（余维 \(1\)）相交的维数至少为 \(\binom{q}{j}-1\)，从而得到 \eqref{eq:pom-multiplicity-composition-replica-explicit-eigs} 的重数下界。
剩余本征值位于 \(S_q\)-对称子空间（维数 \(q+1\)），并由命题 \ref{prop:pom-multiplicity-composition-sq-symmetry-reduction} 的降维给出。证毕。
\end{proof}

% Auto-split: keep each TeX source comfortably below ~600 lines.
%
% This part isolates the (q+1)-dimensional S_q-symmetric exceptional block of the
% q-replica soft-core transfer matrix T^{(q)} (equivalently A^{(q)}), and records
% closed-form Vieta data, asymptotics, and Perron--Doob mixing consequences.
%

\paragraph{例外谱块的 Vieta 刚性与 Perron--Doob 混合常数}

\begin{theorem}[例外 Perron 特征向量在 Hamming 权层上的严格单调性]\label{thm:pom-replica-softcore-exceptional-perron-hamming-monotone}
设 \(A^{(q)}\) 为命题 \ref{prop:pom-multiplicity-composition-sq-symmetry-reduction} 的 \((q+1)\times(q+1)\) 矩阵 \eqref{eq:pom-multiplicity-composition-Aq-def}。
记 \(\rho_q=\rho(A^{(q)})\) 为 Perron 根，并取 Perron 特征向量 \(h^{(q)}=(h_0,\dots,h_q)^{\mathsf T}\succ 0\) 满足
\[
A^{(q)}h^{(q)}=\rho_q h^{(q)}.
\]
则成立严格链
\[
h_0>h_1>\cdots>h_q>0.
\]
\leanverified{paper\_pom\_replica\_softcore\_exceptional\_perron\_hamming\_monotone}
\end{theorem}
\begin{proof}
由特征方程，对任意 \(0\le k\le q-1\) 作差：
\[
\rho_q(h_k-h_{k+1})
=\sum_{\ell=0}^{q}\bigl(A^{(q)}_{k,\ell}-A^{(q)}_{k+1,\ell}\bigr)h_\ell
=\frac12\sum_{\ell=0}^{q}\Bigl(\binom{q-k}{\ell}-\binom{q-k-1}{\ell}\Bigr)h_\ell.
\]
由恒等式 \(\binom{n}{\ell}-\binom{n-1}{\ell}=\binom{n-1}{\ell-1}\)（并约定 \(\binom{n-1}{-1}=0\)），得
\[
\rho_q(h_k-h_{k+1})
=\frac12\sum_{\ell=1}^{q-k}\binom{q-k-1}{\ell-1}\,h_\ell.
\]
右端为严格正数：因 \(h_\ell>0\)，且至少 \(\ell=q-k\) 项的系数为 \(1\)。
故 \(h_k-h_{k+1}>0\) 对所有 \(k\) 成立，从而得到严格单调链。证毕。
\end{proof}

\begin{theorem}[例外谱乘积的完全闭式与三角数符号律]\label{thm:pom-replica-softcore-exceptional-spectrum-product}
沿用命题 \ref{prop:pom-multiplicity-composition-sq-symmetry-reduction} 的
\(S_q\)-对称降维矩阵 \(A^{(q)}\)（定义 \eqref{eq:pom-multiplicity-composition-Aq-def}），
并记其谱（按代数重数计）为
\[
\mathrm{spec}\!\bigl(A^{(q)}\bigr)=\{\lambda_i^{(q)}\}_{i=0}^{q}.
\]
则对任意整数 \(q\ge 2\)，例外谱乘积满足严格恒等式
\begin{equation}\label{eq:pom-replica-softcore-exceptional-product}
\boxed{\;
\prod_{i=0}^{q}\lambda_i^{(q)}
\;=\;
(-1)^{\frac{q(q+1)}{2}}\,2^{-q}. \;}
\end{equation}
等价地，
\[
\det\!\bigl(A^{(q)}\bigr)=(-1)^{\frac{q(q+1)}{2}}\,2^{-q},
\]
从而该乘积与 \(\varphi\) 无关并呈纯二进制刚性。
\leanverified{paper\_pom\_replica\_exceptional\_spectrum\_product\_seeds}
\end{theorem}
\begin{proof}
由 \eqref{eq:pom-multiplicity-composition-Aq-def}，令
\(
\mathbf{1}:=(1,\dots,1)^{\mathsf T}\in\RR^{q+1}
\)
与
\(
b:=(\binom{q}{0},\dots,\binom{q}{q})^{\mathsf T}
\)
并定义矩阵 \(C^{(q)}\in\RR^{(q+1)\times(q+1)}\) 为
\[
C^{(q)}_{k,l}:=\binom{q-k}{l}\qquad(0\le k,l\le q).
\]
则有分解
\begin{equation}\label{eq:pom-replica-softcore-Aq-rankone-plus-C}
A^{(q)}=\frac12\bigl(\mathbf{1}b^{\mathsf T}+C^{(q)}\bigr).
\end{equation}
\par
首先计算 \(\det(C^{(q)})\)。
令 \(R\) 为行反转置换矩阵：\((Rx)_k=x_{q-k}\)。
则
\[
(RC^{(q)})_{k,l}=C^{(q)}_{q-k,l}=\binom{k}{l},
\]
从而 \(RC^{(q)}=P\)，其中 \(P_{k,l}:=\binom{k}{l}\) 为下三角 Pascal 矩阵，故 \(\det(P)=1\)。
而 \(R\) 的行反转需要 \(\frac{q(q+1)}{2}\) 次交换，因此
\begin{equation}\label{eq:pom-replica-softcore-detC}
\det(C^{(q)})=\det(R)\det(P)=(-1)^{\frac{q(q+1)}{2}}.
\end{equation}
\par
其次对 \(\det(\mathbf{1}b^{\mathsf T}+C^{(q)})\) 使用秩一行列式引理：
\[
\det\!\bigl(C^{(q)}+\mathbf{1}b^{\mathsf T}\bigr)=\det(C^{(q)})\Bigl(1+b^{\mathsf T}(C^{(q)})^{-1}\mathbf{1}\Bigr).
\]
注意 \(C^{(q)}e_0=\mathbf{1}\)（其中 \(e_0:=(1,0,\dots,0)^{\mathsf T}\)），因此 \((C^{(q)})^{-1}\mathbf{1}=e_0\)，从而
\(
b^{\mathsf T}(C^{(q)})^{-1}\mathbf{1}=b^{\mathsf T}e_0=\binom{q}{0}=1
\)。
代回得到
\[
\det\!\bigl(C^{(q)}+\mathbf{1}b^{\mathsf T}\bigr)=2\,\det(C^{(q)}).
\]
结合 \eqref{eq:pom-replica-softcore-Aq-rankone-plus-C} 与 \eqref{eq:pom-replica-softcore-detC} 即得
\[
\det(A^{(q)})
=2^{-(q+1)}\det\!\bigl(C^{(q)}+\mathbf{1}b^{\mathsf T}\bigr)
=2^{-q}\det(C^{(q)})
=(-1)^{\frac{q(q+1)}{2}}2^{-q},
\]
等价于 \eqref{eq:pom-replica-softcore-exceptional-product}。证毕。
\end{proof}

\begin{theorem}[例外谱迹的闭式：\texorpdfstring{$2^{q-1}$}{2^{q-1}} 主项与 \texorpdfstring{$\mathbb Z[\varphi]$}{Z[phi]} 校正项]\label{thm:pom-replica-softcore-exceptional-spectrum-trace}
在定理 \ref{thm:pom-replica-softcore-exceptional-spectrum-product} 的设定下，令
\(
\rho_q:=\max\{\lambda_i^{(q)}\}
=\rho\!\bigl(A^{(q)}\bigr)=\rho\!\bigl(T^{(q)}\bigr)
\)
为 Perron 根。
则例外谱和（即例外块迹）满足严格恒等式
\begin{equation}\label{eq:pom-replica-softcore-exceptional-trace}
\boxed{\;
\sum_{i=0}^{q}\lambda_i^{(q)}
\;=\;
2^{q-1}
\;+\;
\frac{\varphi^{q+2}+(-1)^q\varphi^{-q}}{2(\varphi+2)}\; } .
\end{equation}
因此例外块的迹由二项主项 \(2^{q-1}\) 与一个显式的 \(\mathbb Z[\varphi]\)-校正项构成。
\leanverified{paper\_pom\_replica\_exceptional\_spectrum\_trace\_seeds}
\end{theorem}
\begin{proof}
由 \eqref{eq:pom-multiplicity-composition-Aq-def} 直接计算迹：
\[
\Tr\!\bigl(A^{(q)}\bigr)
=\sum_{k=0}^{q}A^{(q)}_{k,k}
=\frac12\sum_{k=0}^{q}\binom{q}{k}
\;+\;\frac12\sum_{k=0}^{q}\binom{q-k}{k}
=2^{q-1}+\frac12\sum_{k=0}^{q}\binom{q-k}{k}.
\]
利用经典恒等式
\(
\sum_{k\ge 0}\binom{q-k}{k}=F_{q+1}
\)
（其中 \(F_n\) 为 Fibonacci 数），得到
\[
\Tr\!\bigl(A^{(q)}\bigr)=2^{q-1}+\frac12F_{q+1}.
\]
再用 Binet 公式 \(F_{q+1}=\frac{\varphi^{q+1}+(-1)^q\varphi^{-(q+1)}}{\sqrt5}\)
与恒等式 \(\varphi+2=\sqrt5\,\varphi\)，可化为
\[
\frac12F_{q+1}
=\frac{\varphi^{q+2}+(-1)^q\varphi^{-q}}{2(\varphi+2)}.
\]
最后由于 \(\Tr(A^{(q)})=\sum_i\lambda_i^{(q)}\)，即得 \eqref{eq:pom-replica-softcore-exceptional-trace}。证毕。
\leanverified{binomial\_row\_sum\_seeds}
\leanverified{alternating\_binomial\_seeds}
\leanverified{triangular\_binomial\_seeds}
\leanverified{paper\_pom\_replica\_trace\_extended}
\end{proof}

\paragraph{二项--几何谱测度的 Fibonacci 塌缩：幂矩、Perron 固定点方程与 Hankel 乘积证书}\label{par:pom-replica-softcore-fibonacci-moment-collapse}
沿用 \( \varphi=\frac{1+\sqrt5}{2} \)、\(\psi=\frac{1-\sqrt5}{2}=-\varphi^{-1}\)。
在命题 \ref{prop:pom-multiplicity-composition-replica-explicit-eigs} 的固定谱簇上，对每个整数 \(q\ge 1\) 引入节点
\begin{equation}\label{eq:pom-replica-softcore-di-nodes}
\boxed{\;
d_i:=\frac12\,\varphi^{q-i}\psi^i=\frac12(-1)^i\varphi^{q-2i}\qquad(0\le i\le q).}
\end{equation}
并定义参数
\begin{equation}\label{eq:pom-replica-softcore-alpha-beta}
\boxed{\;
\alpha:=1+\frac{2}{\sqrt5}=\frac{\varphi^3}{\sqrt5},\qquad
\beta:=1-\frac{2}{\sqrt5}=-\frac{\psi^3}{\sqrt5},\qquad
w_i:=\binom{q}{i}\alpha^{q-i}\beta^i.}
\end{equation}

\begin{proposition}[\texorpdfstring{$S_q$}{Sq}-对称子空间上的张量谱基与二项权重闭式]\label{prop:pom-replica-softcore-sym-eigenbasis-binomial-weights}
\leanverified{paper\_pom\_replica\_softcore\_sym\_eigenbasis\_binomial\_weights}
取整数 \(q\ge 1\)，并令 \(E=\RR^2\)、\(K=\bigl(\begin{smallmatrix}1&1\\[2pt]1&0\end{smallmatrix}\bigr)\)、\(D^{(q)}:=K^{\otimes q}\)。
记 \(u:=(1,1)^{\mathsf T}\in E\) 且 \(\mathbf{1}:=u^{\otimes q}\in E^{\otimes q}\cong \RR^{U_q}\)，并令 \(\mathcal H_{\mathrm{sym}}\subseteq E^{\otimes q}\) 为 \(S_q\) 作用的不变子空间（\(\dim\mathcal H_{\mathrm{sym}}=q+1\)）。
则存在 \(\mathcal H_{\mathrm{sym}}\) 上的正交规范特征基 \(\{f_i\}_{i=0}^{q}\) 使得
\[
D^{(q)}f_i=(2d_i)f_i,
\qquad
d_i=\frac12(-1)^i\varphi^{q-2i},
\]
并且 \(\mathbf{1}\) 在该基下的投影满足
\[
\langle f_i,\mathbf{1}\rangle^2=w_i,
\qquad
\frac{w_i}{2^q}=\binom{q}{i}p^{\,q-i}(1-p)^i,
\qquad
p:=\frac{\alpha}{2}=\frac12+\frac{1}{\sqrt5}.
\]
特别地，\(\sum_{i=0}^q w_i=2^q\)，并且
\[
\prod_{i=0}^{q}w_i
=
\left(\prod_{i=0}^{q}\binom{q}{i}\right)\cdot 5^{-\frac{q(q+1)}{2}}.
\]
\end{proposition}

\begin{proof}
矩阵 \(K\) 的两特征值为 \(\varphi\) 与 \(\psi=-\varphi^{-1}\)，并可取一组正交规范特征向量 \(v_+,v_-\in E\) 使
\(
Kv_+=\varphi v_+
\)、\(
Kv_-=\psi v_-
\)、\(
v_+\perp v_-
\)。
在 \(E^{\otimes q}\) 中，张量积基向量 \(\bigotimes_{j=1}^q v_{\varepsilon_j}\)（\(\varepsilon_j\in\{+,-\}\)）为 \(D^{(q)}=K^{\otimes q}\) 的本征向量，其本征值为 \(\varphi^{q-i}\psi^i\)，其中 \(i\) 为符号 “\(-\)” 的出现次数。
将这些向量对 \(S_q\) 的置换作用作对称化并按 \(i\) 分层，即可在 \(\mathcal H_{\mathrm{sym}}\) 上得到正交规范特征基 \(\{f_i\}_{i=0}^{q}\)，其本征值为 \(\varphi^{q-i}\psi^i=(-1)^i\varphi^{q-2i}=2d_i\)。
\par
另一方面，\(\mathbf{1}=u^{\otimes q}\)。
将 \(u\) 在正交基 \(\{v_+,v_-\}\) 下展开为 \(u=\alpha_0 v_+ +\beta_0 v_-\)，则
\(
u^{\otimes q}
=\sum_{i=0}^q \sqrt{\binom{q}{i}}\alpha_0^{\,q-i}\beta_0^{\,i} f_i
\)
且
\(
w_i=\langle f_i,\mathbf{1}\rangle^2=\binom{q}{i}\alpha_0^{\,2(q-i)}\beta_0^{\,2i}
\)。
在当前归一化下有
\(
\alpha_0^{2}=\alpha
\)、\(
\beta_0^{2}=\beta
\)（见 \eqref{eq:pom-replica-softcore-alpha-beta}），从而 \(w_i=\binom{q}{i}\alpha^{q-i}\beta^i\)，与 \eqref{eq:pom-replica-softcore-alpha-beta} 一致。
由 \(\alpha+\beta=2\) 得 \(\sum_i w_i=(\alpha+\beta)^q=2^q\)。
又由 \(\alpha\beta=1/5\) 得
\[
\prod_{i=0}^{q}w_i
=\left(\prod_{i=0}^{q}\binom{q}{i}\right)\alpha^{\sum_{i}(q-i)}\beta^{\sum_i i}
=\left(\prod_{i=0}^{q}\binom{q}{i}\right)(\alpha\beta)^{\frac{q(q+1)}{2}}
=\left(\prod_{i=0}^{q}\binom{q}{i}\right)5^{-\frac{q(q+1)}{2}}.
\]
最后，令 \(p=\alpha/2\) 即得 \(w_i/2^q=\binom{q}{i}p^{q-i}(1-p)^i\)。证毕。
\end{proof}

\begin{proposition}[秩一扰动的割线方程接口（例外谱的一维刻画）]\label{prop:pom-replica-softcore-exceptional-secant-equation}
\leanverified{paper\_pom\_replica\_softcore\_exceptional\_secant\_equation}
取整数 \(q\ge 2\)，并令 \(\rho_q=\rho(A^{(q)})\) 为例外块的 Perron 根（见定理 \ref{thm:pom-replica-softcore-exceptional-spectrum-trace} 的记号）。
则 \(\rho_q\) 满足严格割线方程
\begin{equation}\label{eq:pom-replica-softcore-secant-equation}
\boxed{\;
\sum_{i=0}^{q}\frac{w_i}{\rho_q-d_i}=2.}
\end{equation}
\end{proposition}
\begin{proof}[证明要点]
由命题 \ref{prop:pom-multiplicity-composition-replica-explicit-eigs}，固定谱簇在 \((S_q)\)-对称子空间上以 \(\{d_i\}_{i=0}^{q}\) 为节点呈对角化；
而 \(T^{(q)}=\frac12(\mathbf{1}\mathbf{1}^{\mathsf T}+D^{(q)})\) 的 \(\mathbf{1}\mathbf{1}^{\mathsf T}\) 项在该子空间上给出秩一耦合。
在该谱基下应用秩一行列式引理（等价于 rank-one resolvent 恒等式），得到例外特征值 \(\lambda\) 满足
\(
1-\frac12\sum_i \frac{w_i}{\lambda-d_i}=0
\)，即 \eqref{eq:pom-replica-softcore-secant-equation}。
权重 \(w_i\) 的闭式由对称张量本征基与 \(\mathbf{1}\) 的投影系数显式计算得到，见 \eqref{eq:pom-replica-softcore-alpha-beta}。证毕。
\end{proof}

\begin{corollary}[割线方程的期望形式]\label{cor:pom-replica-softcore-secant-equation-expectation-form}
\leanverified{paper\_pom\_replica\_softcore\_secant\_equation\_expectation\_form}
取 \(q\ge 2\)，令概率权重 \(\pi_i:=w_i/2^q\) 并令随机变量 \(I\sim \mathrm{Bin}(q,1-p)\)（其中 \(p=\frac12+\frac{1}{\sqrt5}\)）。
定义几何节点随机变量 \(X:=d_I\)（节点 \(d_i\) 见 \eqref{eq:pom-replica-softcore-di-nodes}）。
则例外块 Perron 根 \(\rho_q\) 的割线方程 \eqref{eq:pom-replica-softcore-secant-equation} 等价于标量期望方程
\[
\boxed{\;
\mathbb E\!\left[\frac{1}{\rho_q-X}\right]=2^{1-q}.}
\]
\end{corollary}

\begin{proof}
由命题 \ref{prop:pom-replica-softcore-exceptional-secant-equation}，
\(
\sum_i \frac{w_i}{\rho_q-d_i}=2
\)。
两边同除以 \(2^q\) 并用 \(\pi_i=w_i/2^q\) 得
\(
\sum_i \frac{\pi_i}{\rho_q-d_i}=2^{1-q}
\)。
由 \(\pi_i=\binom{q}{i}p^{q-i}(1-p)^i\)（命题 \ref{prop:pom-replica-softcore-sym-eigenbasis-binomial-weights}）可识别 \(\pi_i\) 为 \(I\sim\mathrm{Bin}(q,1-p)\) 的质量函数，故左端为 \(\mathbb E[(\rho_q-d_I)^{-1}]\)，即所示期望形式。证毕。
\end{proof}

\begin{theorem}[二项--几何谱测度的 Fibonacci 幂矩塌缩]\label{thm:pom-replica-softcore-fibonacci-power-moment-collapse}
\leanverified{paper\_pom\_replica\_softcore\_fibonacci\_power\_moment\_collapse}
在 \eqref{eq:pom-replica-softcore-di-nodes}--\eqref{eq:pom-replica-softcore-alpha-beta} 的记号下，对任意整数 \(m\ge 0\) 成立严格恒等式
\begin{equation}\label{eq:pom-replica-softcore-moment-collapse}
\boxed{\;
\sum_{i=0}^{q} w_i\,d_i^{\,m}
\,=\,2^{-m}\,F_{m+3}^{\,q}.}
\end{equation}
因此该谱测度的全部幂矩都落在 \(\QQ\) 中，右端不含 \(\sqrt5\) 的代数延拓项。
\end{theorem}
\begin{proof}
由 \(\psi=-\varphi^{-1}\) 得 \(d_i^{\,m}=2^{-m}\varphi^{m(q-i)}\psi^{mi}\)。
于是
\[
\sum_{i=0}^{q} w_i\,d_i^{\,m}
=2^{-m}\sum_{i=0}^{q}\binom{q}{i}(\alpha\varphi^m)^{q-i}(\beta\psi^m)^i
=2^{-m}(\alpha\varphi^m+\beta\psi^m)^{q}.
\]
由 \eqref{eq:pom-replica-softcore-alpha-beta} 的 \(\alpha=\varphi^3/\sqrt5\)、\(\beta=-\psi^3/\sqrt5\)，有
\[
\alpha\varphi^m+\beta\psi^m=\frac{\varphi^{m+3}-\psi^{m+3}}{\sqrt5}=F_{m+3},
\]
代回即得 \eqref{eq:pom-replica-softcore-moment-collapse}。证毕。
\end{proof}

\begin{theorem}[Perron 特征值的纯 Fibonacci 固定点方程]\label{thm:pom-replica-softcore-perron-fibonacci-fixed-point}
\leanverified{paper\_pom\_replica\_softcore\_perron\_fibonacci\_fixed\_point}
取整数 \(q\ge 2\)，令 \(\rho_q=\rho(A^{(q)})\) 为例外块 Perron 根。
则存在唯一 \(x_q>1\) 使得
\begin{equation}\label{eq:pom-replica-softcore-rho-xq}
\boxed{\;\rho_q=2^{q-1}x_q\;}
\end{equation}
并且 \(x_q\) 满足严格收敛的固定点方程
\begin{equation}\label{eq:pom-replica-softcore-xq-fixed-point}
\boxed{\;
x_q
=1+\sum_{m=1}^{\infty}\left(\frac{F_{m+3}}{2^{m+1}}\right)^{q}x_q^{-m}.}
\end{equation}
\end{theorem}
\begin{proof}
由命题 \ref{prop:pom-multiplicity-composition-sq-symmetry-reduction}，有 \(\rho_q=\rho(T^{(q)})\)。
取 \(\mathbf{1}\in\RR^{U_q}\) 为全 \(1\) 向量作 Rayleigh 商，得
\[
\rho_q
\ge \frac{\mathbf 1^{\mathsf T}T^{(q)}\mathbf 1}{\mathbf 1^{\mathsf T}\mathbf 1}
=\frac{4^q+3^q}{2^{q+1}}
=2^{q-1}+\frac{3^q}{2^{q+1}}
>\frac{\varphi^q}{2}
=\max_{0\le i\le q} d_i.
\]
因此对每个 \(i\) 可作几何级数展开
\[
\frac{1}{\rho_q-d_i}=\frac{1}{\rho_q}\sum_{m=0}^{\infty}\left(\frac{d_i}{\rho_q}\right)^{m},
\qquad \text{且一致绝对收敛。}
\]
将之代入 \eqref{eq:pom-replica-softcore-secant-equation} 并交换求和次序得
\[
2=\frac{1}{\rho_q}\sum_{m=0}^{\infty}\rho_q^{-m}\sum_{i=0}^{q}w_i d_i^{\,m}.
\]
由定理 \ref{thm:pom-replica-softcore-fibonacci-power-moment-collapse}，
\(
\sum_i w_i d_i^{\,m}=2^{-m}F_{m+3}^q
\)，
从而
\[
2=\frac{1}{\rho_q}\sum_{m=0}^{\infty}\left(\frac{F_{m+3}}{2}\right)^{q}\left(\frac{1}{2\rho_q}\right)^m.
\]
记 \(\rho_q=2^{q-1}x_q\)，整理并移去 \(m=0\) 项（注意 \(F_3=2\)）即得 \eqref{eq:pom-replica-softcore-xq-fixed-point}。
右端关于 \(x\in(1,\infty)\) 严格递减而左端为严格递增函数 \(x\mapsto x\)，故正解唯一。证毕。
\end{proof}

\begin{corollary}[Perron 根的两项主导展开与显式余项界]\label{cor:pom-replica-softcore-rho-two-term-remainder}
\leanverified{paper\_pom\_replica\_softcore\_rho\_two\_term\_remainder}
定义
\[
b_m:=\frac{F_{m+3}}{2^{m+1}}\quad(m\ge 1),
\qquad
\varepsilon_q:=\sum_{m=1}^{\infty}b_m^{\,q}.
\]
则 \(x_q\in(1,1+\varepsilon_q)\)，并且存在余项 \(R_q\) 使
\begin{equation}\label{eq:pom-replica-softcore-rho-two-term}
\rho_q
=2^{q-1}+\frac{3^q}{2^{q+1}}+R_q,
\end{equation}
其中余项满足显式上界
\begin{equation}\label{eq:pom-replica-softcore-rho-remainder-bound}
\boxed{\;
0<R_q
<2^{q-1}\sum_{m=2}^{\infty}b_m^{\,q}
\le \frac{5^q}{2^{2q+1}}\cdot\frac{1}{1-(\varphi/2)^q}.}
\end{equation}
因此当 \(q\to\infty\) 时有严格渐近
\[
\rho_q
=2^{q-1}+\frac{3^q}{2^{q+1}}+O\!\left(\frac{5^q}{2^{2q}}\right).
\]
\end{corollary}
\begin{proof}
由定理 \ref{thm:pom-replica-softcore-perron-fibonacci-fixed-point} 的固定点方程，
\(
x_q-1=\sum_{m\ge 1}b_m^{\,q}x_q^{-m}
\)。
因 \(x_q>1\)，得 \(0<x_q-1<\sum_{m\ge 1}b_m^{\,q}=\varepsilon_q\)，从而 \(x_q\in(1,1+\varepsilon_q)\)。
再将首项 \(m=1\) 分离并乘以 \(2^{q-1}\)，用 \(b_1=F_4/2^2=3/4\) 得到 \eqref{eq:pom-replica-softcore-rho-two-term}，并且 \(R_q>0\)。
对余项上界，用 \(x_q^{-m}\le 1\) 得
\(
R_q<2^{q-1}\sum_{m\ge 2}b_m^{\,q}
\)。
又由 Fibonacci 比值单调性 \(F_{n+1}/F_n<\varphi\)（\(n\ge 1\)）推出
\(
b_{m+1}/b_m<\varphi/2
\)，
故 \(\sum_{m\ge 2}b_m^{\,q}\) 被等比级数控制并得到 \eqref{eq:pom-replica-softcore-rho-remainder-bound}。证毕。
\end{proof}

\begin{theorem}[Perron 根的 Schur--Feshbach 方程与组合--矩级数]\label{thm:pom-replica-softcore-perron-feshbach-composition-moment-series}
\leanverified{paper\_pom\_replica\_softcore\_perron\_feshbach\_composition\_moment\_series}
令 \(n:=2^q\)，并在 \(\RR^{U_q}\) 上记 \(u:=n^{-1/2}\mathbf{1}\)、\(P:=uu^{\mathsf T}\)、\(Q:=I-P\)。
令 \(D^{(q)}=K^{\otimes q}\) 且 \(E:=\frac12D^{(q)}\)，则
\[
T^{(q)}=\frac12\bigl(\mathbf{1}\mathbf{1}^{\mathsf T}+D^{(q)}\bigr)=2^{q-1}P+E.
\]
记 \(\lambda_0^{(q)}:=\rho(T^{(q)})=\rho_q\)。
定义幂矩
\[
m_k(q):=\langle u,E^k u\rangle=\frac{1}{2^{q+k}}\,F_{k+3}^{\,q}\qquad(k\ge 1),
\]
并对 \(m\ge 1\) 定义组合系数
\[
\kappa_{m}(q):=\sum_{r=1}^{m}(-1)^{r-1}\!\!\sum_{\substack{\ell_1+\cdots+\ell_r=m\\ \ell_i\ge 1}}\ \prod_{j=1}^{r} m_{\ell_j}(q),
\]
其中内层求和遍历将 \(m\) 写成正整数和的有序组合（composition）。
则 \(\lambda_0^{(q)}\) 满足严格标量方程
\[
\boxed{\;
\lambda_0^{(q)}
=2^{q-1}+m_1(q)+\sum_{s=0}^{\infty}\frac{\kappa_{s+2}(q)}{\bigl(\lambda_0^{(q)}\bigr)^{s+1}}\;}
\]
且该级数绝对收敛。
其前三个非平凡系数可写为
\[
\kappa_2(q)=m_2(q)-m_1(q)^2,\qquad
\kappa_3(q)=m_3(q)-2m_1(q)m_2(q)+m_1(q)^3,
\]
\[
\kappa_4(q)=m_4(q)-m_2(q)^2-2m_1(q)m_3(q)+3m_1(q)^2m_2(q)-m_1(q)^4,
\]
其中
\(
m_1(q)=\frac{3^q}{2^{q+1}}
\)、\(
m_2(q)=\frac{5^q}{2^{q+2}}
\)、\(
m_3(q)=\frac{8^q}{2^{q+3}}
\)、\(
m_4(q)=\frac{13^q}{2^{q+4}}
\)。
\end{theorem}

\begin{proof}
将 \(\RR^{U_q}=\mathrm{span}\{u\}\oplus u^\perp\) 分解，并在该分解下写
\[
T^{(q)}=
\begin{pmatrix}
2^{q-1}+m_1(q) & b^{\mathsf T}\\
b & C
\end{pmatrix},
\qquad
b:=QEu,\qquad C:=QEQ.
\]
对特征值 \(\lambda_0^{(q)}\) 取 Schur 补得到
\[
\lambda_0^{(q)}=2^{q-1}+m_1(q)+b^{\mathsf T}\bigl(\lambda_0^{(q)}I-C\bigr)^{-1}b.
\]
由 \(\|C\|\le \|E\|=\frac12\|D^{(q)}\|=\frac12\varphi^q\) 且 \(\lambda_0^{(q)}\ge 2^{q-1}\)，得
\(
\|C\|/\lambda_0^{(q)}\le (\varphi/2)^q<1
\)，因此
\[
(\lambda_0^{(q)}I-C)^{-1}
=\sum_{s=0}^{\infty}\frac{C^s}{(\lambda_0^{(q)})^{s+1}}
\qquad\text{（算子范数绝对收敛）}.
\]
代入得到
\[
\lambda_0^{(q)}=2^{q-1}+m_1(q)+\sum_{s=0}^{\infty}\frac{b^{\mathsf T}C^s b}{(\lambda_0^{(q)})^{s+1}}.
\]
最后，利用 \(Q=I-P\) 的展开与恒等式
\(
\langle u,XPYu\rangle=\langle u,Xu\rangle\langle u,Yu\rangle
\)，可将 \(b^{\mathsf T}C^s b=\langle u,EQ(QEQ)^sQE u\rangle\) 展开为“把 \(s+2\) 个 \(E\) 依次分割成若干块”的有序组合求和，
其符号为 \((-1)^{r-1}\)，并贡献块积 \(\prod_j \langle u,E^{\ell_j}u\rangle=\prod_j m_{\ell_j}(q)\)。
因此 \(b^{\mathsf T}C^s b=\kappa_{s+2}(q)\)，从而得到所示级数方程。
\par
幂矩闭式由 \(E^k=2^{-k}(K^k)^{\otimes q}\) 与 \(u=n^{-1/2}\mathbf{1}\) 给出：
\(
m_k(q)=2^{-q-k}\bigl(\mathbf{1}^{\mathsf T}K^k\mathbf{1}\bigr)^q
\)，而 \(\mathbf{1}^{\mathsf T}K^k\mathbf{1}=F_{k+3}\)（由 \(K^k\) 的 Fibonacci 形式）即得 \(m_k(q)=2^{-(q+k)}F_{k+3}^q\)。
系数 \(\kappa_2,\kappa_3,\kappa_4\) 由组合定义直接展开即得。证毕。
\end{proof}

\begin{proposition}[Perron 特征向量的均匀化定量界]\label{prop:pom-replica-softcore-perron-eigenvector-uniformization-angle-bound}
设 \(v^{(q)}\) 为 \(T^{(q)}\) 的 Perron 特征向量，取 \(\|v^{(q)}\|_2=1\) 且分量全正，并仍记 \(u=2^{-q/2}\mathbf{1}\)。
则有定量界
\[
\sin\angle\bigl(v^{(q)},u\bigr)
\le
\frac{\|QEu\|_2}{2^{q-1}-\|E\|}
=
\frac{\sqrt{\frac{5^q}{2^{q+2}}-\frac{9^q}{2^{2q+2}}}}{2^{q-1}-\frac12\varphi^q},
\]
从而
\[
\|v^{(q)}-u\|_2\le \sqrt2\,\sin\angle\bigl(v^{(q)},u\bigr)
\le
\frac{\sqrt{\frac{5^q}{2^{q+1}}-\frac{9^q}{2^{2q+1}}}}{2^{q-1}-\frac12\varphi^q}.
\]
\leanverified{paper\_pom\_replica\_softcore\_perron\_eigenvector\_uniformization\_angle\_bound}
\end{proposition}

\begin{proof}
由 \(T^{(q)}=2^{q-1}P+E\) 与 \(T^{(q)}v^{(q)}=\lambda_0^{(q)}v^{(q)}\)，对等式取 \(Q\) 投影得
\[
(\lambda_0^{(q)}I-QEQ)\,Qv^{(q)}=QEu\cdot \langle u,v^{(q)}\rangle.
\]
因此
\[
\|Qv^{(q)}\|_2\le \|(\lambda_0^{(q)}I-QEQ)^{-1}\|\cdot \|QEu\|_2
\le \frac{\|QEu\|_2}{\lambda_0^{(q)}-\|E\|}
\le \frac{\|QEu\|_2}{2^{q-1}-\|E\|},
\]
其中使用了 \(\|QEQ\|\le \|E\|\) 与 \(\lambda_0^{(q)}\ge 2^{q-1}\)。
又 \(\|Qv^{(q)}\|_2=\sin\angle(v^{(q)},u)\)，得第一式。
\par
对 \(\|QEu\|_2\)，注意
\(
\|Eu\|_2^2=\langle u,E^2u\rangle=m_2(q)=\frac{5^q}{2^{q+2}}
\) 且 \(\langle u,Eu\rangle=m_1(q)=\frac{3^q}{2^{q+1}}\)，从而
\[
\|QEu\|_2^2=\|Eu\|_2^2-\langle u,Eu\rangle^2
=\frac{5^q}{2^{q+2}}-\frac{9^q}{2^{2q+2}}.
\]
最后用 \(\|E\|=\frac12\|D^{(q)}\|=\frac12\varphi^q\) 即得所示显式界；并由 \(\|v-u\|_2\le \sqrt2\,\sin\angle(v,u)\) 得第二式。证毕。
\end{proof}

\begin{theorem}[Fibonacci 幂 Hankel 行列式的完全乘积公式]\label{thm:pom-replica-softcore-fibonacci-power-hankel-determinant}
\leanverified{paper\_pom\_replica\_softcore\_fibonacci\_power\_hankel\_determinant}
对任意整数 \(q\ge 1\)，定义 \((q+1)\times(q+1)\) Hankel 矩阵
\[
H_q:=\bigl(F_{i+j+3}^{\,q}\bigr)_{0\le i,j\le q}.
\]
则成立精确闭式
\begin{equation}\label{eq:pom-replica-softcore-fibpower-hankel-det}
\boxed{\;
\det(H_q)
=\left(\prod_{i=0}^{q}\binom{q}{i}\right)
\left(\prod_{k=1}^{q}F_k^{\,2(q+1-k)}\right).}
\end{equation}
右端为整数的纯乘积表达，不含任何代数数域延拓。
\end{theorem}
\begin{proof}
由 Binet 公式与 \eqref{eq:pom-replica-softcore-alpha-beta} 的 \(\alpha,\beta\) 可写
\(
F_{n+3}=\alpha\varphi^n+\beta\psi^n
\)。
从而对每个整数 \(n\ge 0\) 有指数和表示
\[
F_{n+3}^{\,q}
=\sum_{i=0}^{q}\binom{q}{i}\alpha^{q-i}\beta^i(\varphi^{q-i}\psi^i)^{n}
=\sum_{i=0}^{q}c_i r_i^{\,n},
\]
其中 \(r_i:=\varphi^{q-i}\psi^i\)、\(c_i:=\binom{q}{i}\alpha^{q-i}\beta^i\)。
令 \(V=(r_i^{\,j})_{0\le i,j\le q}\) 且 \(C=\mathrm{diag}(c_0,\dots,c_q)\)，则 \(H_q=V^{\mathsf T}CV\)，因而
\[
\det(H_q)=(\det V)^2\prod_{i=0}^{q}c_i
=\left(\prod_{0\le i<j\le q}(r_j-r_i)^2\right)\left(\prod_{i=0}^{q}c_i\right).
\]
计算 \(\prod_i c_i=(\prod_i\binom{q}{i})(\alpha\beta)^{q(q+1)/2}\) 且 \(\alpha\beta=1/5\)，得
\(
\prod_i c_i=(\prod_i\binom{q}{i})\,5^{-q(q+1)/2}
\)。
另一方面，把 \(r_i=\varphi^q(\psi/\varphi)^i\) 写为几何节点并用 Vandermonde 分解，
再用恒等式 \(1-(-\varphi^{-2})^k=\sqrt5\,F_k/\varphi^{k}\) 整理，可得
\(
\prod_{i<j}(r_j-r_i)^2
=5^{q(q+1)/2}\prod_{k=1}^{q}F_k^{2(q+1-k)}
\)；
两者的 \(5\)-因子精确抵消，得到 \eqref{eq:pom-replica-softcore-fibpower-hankel-det}。证毕。
\end{proof}

\begin{corollary}[交错几何点的 Vandermonde 平方塌缩为纯 Fibonacci 乘积]\label{cor:pom-replica-softcore-alternating-vandermonde-fibonacci-product}
\leanverified{paper\_pom\_replica\_softcore\_alternating\_vandermonde\_fibonacci\_product}
令
\(
\xi_i:=(-1)^i\varphi^{q-2i}\ (0\le i\le q)
\)。
则成立严格恒等式
\begin{equation}\label{eq:pom-replica-softcore-alternating-vandermonde}
\boxed{\;
\prod_{0\le i<j\le q}(\xi_j-\xi_i)^2
=5^{q(q+1)/2}\prod_{k=1}^{q}F_k^{2(q+1-k)}\in\ZZ.}
\end{equation}
\end{corollary}
\begin{proof}
注意 \(r_i=\varphi^{q-i}\psi^i=\xi_i\)（由 \(\psi^i=(-1)^i\varphi^{-i}\)），
故定理 \ref{thm:pom-replica-softcore-fibonacci-power-hankel-determinant} 证明中的 Vandermonde 平方闭式直接给出 \eqref{eq:pom-replica-softcore-alternating-vandermonde}。证毕。
\end{proof}

\begin{proposition}[节点多项式判别式的域范数：Fibonacci 指数分解]\label{prop:pom-replica-softcore-node-discriminant-field-norm}
令 \(K=\QQ(\sqrt5)\)，并定义节点多项式
\[
G_q(x):=\prod_{i=0}^{q}\bigl(x-\xi_i\bigr)\in K[x],
\qquad
\xi_i:=(-1)^i\varphi^{q-2i}.
\]
则其判别式 \(\mathrm{Disc}(G_q)\in K\) 的域范数满足
\[
\boxed{\;
\mathrm N_{K/\QQ}\!\bigl(\mathrm{Disc}(G_q)\bigr)
=5^{q(q+1)}\prod_{d=1}^{q}F_d^{\,4(q+1-d)}\; }.
\]
\end{proposition}

\begin{proof}
由判别式定义，
\(
\mathrm{Disc}(G_q)=\prod_{0\le i<j\le q}(\xi_j-\xi_i)^2
\)。
推论 \ref{cor:pom-replica-softcore-alternating-vandermonde-fibonacci-product} 给出该乘积在 \(K\) 中实际上属于 \(\ZZ\) 且等于
\(
5^{q(q+1)/2}\prod_{d=1}^{q}F_d^{2(q+1-d)}
\)。
因此 \(\mathrm{Disc}(G_q)\in\ZZ\subset K\)，从而
\(
\mathrm N_{K/\QQ}(\mathrm{Disc}(G_q))=\mathrm{Disc}(G_q)^2
\)，即得到所示公式。证毕。
\end{proof}

\begin{proposition}[例外特征多项式与节点多项式的结果式：域范数的完全因子化]\label{prop:pom-replica-softcore-resultant-factorization-field-norm}
令 \(K=\QQ(\sqrt5)\)，并保持节点多项式 \(G_q\) 的定义。
记 \(D^{(q)}=K^{\otimes q}\) 在 \(\mathcal H_{\mathrm{sym}}\) 上的谱为 \(\{2d_i\}_{i=0}^q\)，并令
\[
P_q(x):=\det\!\bigl(xI-(J^{(q)}+D^{(q)})|_{\mathcal H_{\mathrm{sym}}}\bigr)\in K[x],
\qquad
J^{(q)}:=\mathbf{1}\mathbf{1}^{\mathsf T}.
\]
则成立域范数完全因子化：
\[
\boxed{\;
\mathrm N_{K/\QQ}\!\bigl(\mathrm{Res}(P_q,G_q)\bigr)
=\left(\prod_{i=0}^{q}\binom{q}{i}^2\right)\cdot
\left(\prod_{d=1}^{q}F_d^{\,4(q+1-d)}\right)\; }.
\]
\end{proposition}

\begin{proof}
在命题 \ref{prop:pom-replica-softcore-sym-eigenbasis-binomial-weights} 的正交基 \(\{f_i\}\) 下，
\(
D^{(q)}|_{\mathcal H_{\mathrm{sym}}}=\mathrm{diag}(2d_0,\dots,2d_q)
\) 且 \(J^{(q)}|_{\mathcal H_{\mathrm{sym}}}=gg^{\mathsf T}\) 其中 \(g_i=\langle f_i,\mathbf{1}\rangle\)、\(g_i^2=w_i\)。
由秩一行列式引理，
\[
P_q(x)=\prod_{i=0}^{q}(x-2d_i)\left(1-\sum_{i=0}^{q}\frac{w_i}{x-2d_i}\right).
\]
因此对每个 \(i\)，在 \(x=2d_i\) 处取极限得到
\[
P_q(2d_i)=-w_i\prod_{j\ne i}(2d_i-2d_j).
\]
从而
\[
\mathrm{Res}(P_q,G_q)
=\prod_{i=0}^{q}P_q(\xi_i)
=\pm\left(\prod_{i=0}^{q}w_i\right)\cdot \mathrm{Disc}(G_q),
\]
其中 \(\xi_i=2d_i\)，而符号对域范数无影响。
对域范数，用 \(\mathrm N_{K/\QQ}(w_i)=\binom{q}{i}^2\cdot 5^{-q}\)（由命题 \ref{prop:pom-replica-softcore-sym-eigenbasis-binomial-weights} 的 \(\alpha\beta=1/5\) 与共轭交换）得
\[
\mathrm N_{K/\QQ}\!\left(\prod_{i=0}^q w_i\right)
=\frac{\prod_{i=0}^q\binom{q}{i}^2}{5^{q(q+1)}}.
\]
再由命题 \ref{prop:pom-replica-softcore-node-discriminant-field-norm} 代入
\(
\mathrm N_{K/\QQ}(\mathrm{Disc}(G_q))=5^{q(q+1)}\prod_{d=1}^q F_d^{4(q+1-d)}
\)，相乘即得结论。证毕。
\end{proof}

\begin{remark}[审计脚本]\label{rem:pom-replica-softcore-fibonacci-moment-collapse-audit-script}
式 \eqref{eq:pom-replica-softcore-moment-collapse}、\eqref{eq:pom-replica-softcore-xq-fixed-point}、\eqref{eq:pom-replica-softcore-fibpower-hankel-det} 与 \eqref{eq:pom-replica-softcore-alternating-vandermonde} 的一致性可由脚本
\texttt{scripts/exp\_pom\_replica\_softcore\_fibonacci\_moment\_collapse\_audit.py}
在小 \(q\) 上进行严格核验。
\end{remark}

\endinput


\begin{theorem}[第二例外特征值的主阶渐近：向 \texorpdfstring{$\varphi^q/2$}{phi^q/2} 的指数贴近]\label{thm:pom-replica-softcore-second-exceptional-eigenvalue-asymptotics}
令 \(\nu_q\) 为例外谱中次大特征值（即 \((S_q)\)-对称子空间上仅次于 \(\rho_q\) 的特征值）。
则 \(0<\nu_q<\varphi^q/2\)，并且存在严格渐近展开
\begin{equation}\label{eq:pom-replica-softcore-nuq-asymptotics}
\nu_q
=\frac{\varphi^{q}}{2}\Bigg(1-r^{q}
+O\!\bigl(q\,r^{2q}\bigr)\Bigg),
\qquad q\to\infty,
\end{equation}
其中
\(
r:=\frac{5+2\sqrt5}{10}
=\frac12\bigl(1+\frac{2\sqrt5}{5}\bigr)\in(0,1)
\)。
换言之，\(\nu_q\) 在相对误差尺度上以底数 \(r\) 指数趋近于 \(\varphi^q/2\)。
\end{theorem}
\begin{proof}[证明要点]
由命题 \ref{prop:pom-multiplicity-composition-replica-explicit-eigs}，
\(
T^{(q)}=\frac12(\mathbf{1}\mathbf{1}^{\mathsf T}+K^{\otimes q})
\)
且 \(K\) 可在本征基 \(\{v^+,v^-\}\) 上对角化（定理 \ref{thm:pom-replica-softcore-tensor-eigenbasis}）。
将 \(S_q\)-对称子空间在该张量谱层上坐标化后，例外块可视为对 \(K^{\otimes q}\) 最大正谱层的秩一耦合扰动；
在 \(q\to\infty\) 下，主耦合系数由单坐标重叠比决定，从而出现底数 \(r\) 的指数校正项，
二阶误差由高阶谱层的指数抑制与标准 resolvent 展开给出。证毕。
\end{proof}

\begin{theorem}[谱隙比率的指数律与对数分离速率]\label{thm:pom-replica-softcore-gap-ratio-log-separation}
令 \(\rho_q\) 与 \(\nu_q\) 如上，则
\begin{equation}\label{eq:pom-replica-softcore-gap-ratio}
\frac{\nu_q}{\rho_q}
=\Big(\frac{\varphi}{2}\Big)^{q}\Big(1+o(1)\Big),
\qquad q\to\infty,
\end{equation}
并且对数分离满足
\begin{equation}\label{eq:pom-replica-softcore-gap-log-separation}
\log\frac{\rho_q}{\nu_q}
= q\log\frac{2}{\varphi}+o(q).
\end{equation}
因此例外谱内的主次两特征值在对数尺度上以速率 \(\log(2/\varphi)\) 线性分离。
\end{theorem}
\begin{proof}
由定理 \ref{thm:pom-replica-softcore-second-exceptional-eigenvalue-asymptotics}，
\(
\nu_q=\frac{\varphi^q}{2}(1+o(1))
\)。
另一方面，\(T^{(q)}\) 为对称矩阵且满足
\(
T^{(q)}=\frac12(\mathbf{1}\mathbf{1}^{\mathsf T}+K^{\otimes q})
\)
（命题 \ref{prop:pom-multiplicity-composition-replica-explicit-eigs}）。
由 Weyl 扰动界（对称矩阵的特征值稳定性），有
\[
\bigl|\rho_q-2^{q-1}\bigr|
\le \frac12\,\rho\!\bigl(K^{\otimes q}\bigr)
=\frac{\varphi^q}{2},
\]
从而 \(\rho_q=2^{q-1}(1+o(1))\)（因 \((\varphi/2)^q\to 0\)）。
将两式相除即得 \eqref{eq:pom-replica-softcore-gap-ratio}，
再取对数得到 \eqref{eq:pom-replica-softcore-gap-log-separation}。证毕。
\end{proof}

\begin{theorem}[Perron 方向的定量稳定性：与均匀向量的夹角指数小]\label{thm:pom-replica-softcore-perron-angle-uniform}
令 \(N=2^q\)，并记单位均匀向量
\[
u:=\frac{1}{\sqrt{N}}\mathbf 1\in\RR^{U_q}.
\]
令 \(v\) 为 \(T^{(q)}\) 的单位 Perron 特征向量（取 \(\|v\|_2=1\)、\(v\succ 0\)）。
则对任意 \(q\ge 1\) 有角度估计
\[
\sin\angle(u,v)\ \le\ \Bigl(\frac{\varphi}{2}\Bigr)^{q},
\qquad \bigl(\varphi=\tfrac{1+\sqrt5}{2}\bigr).
\]
从而
\[
|u^{\mathsf T}v|\ \ge\ \sqrt{1-\Bigl(\frac{\varphi}{2}\Bigr)^{2q}}.
\]
\leanverified{paper\_pom\_replica\_softcore\_perron\_angle\_uniform}
\end{theorem}
\begin{proof}
写
\[
T^{(q)}=\frac12J+\frac12D,\qquad
J:=\mathbf 1\mathbf 1^{\mathsf T},\qquad D:=K^{\otimes q}.
\]
矩阵 \(\frac12J\) 的谱为 \(\{2^{q-1},0,\dots,0\}\)，其 Perron 特征向量为 \(u\)，并且 \(\frac12J\) 的 Perron 特征值与其余谱的间隙为 \(2^{q-1}\)。
另一方面，由张量谱乘法性，\(\|D\|_2=\rho(D)=\varphi^q\)，故
\(
\|D/2\|_2=\varphi^q/2
\)。
对称谱扰动的 \(\sin\Theta\) 型估计（Davis--Kahan）给出
\[
\sin\angle(u,v)\le \frac{\|D/2\|_2}{2^{q-1}}=\Bigl(\frac{\varphi}{2}\Bigr)^q.
\]
再由 \(|u^{\mathsf T}v|=\cos\angle(u,v)\) 得到第二式。证毕。
\end{proof}

\begin{theorem}[Perron--Doob 变换的可逆马尔可夫核与相关衰减常数]\label{thm:pom-replica-softcore-doob-reversible-markov-corr}
令 \(\psi^{(q)}\succ 0\) 为 \(T^{(q)}\) 的 Perron 特征向量（取 \(\|\psi^{(q)}\|_2=1\)）。
定义 Perron--Doob 变换核
\[
\mathcal{P}^{(q)}(u,v)
:=\frac{T^{(q)}_{u,v}\,\psi^{(q)}_{v}}{\rho_q\,\psi^{(q)}_{u}},
\qquad u,v\in U_q.
\]
则 \(\mathcal{P}^{(q)}\) 为可逆马尔可夫核，其平稳分布为
\[
\pi^{(q)}(u)\propto \bigl(\psi^{(q)}_{u}\bigr)^2.
\]
并且对任意平方可积可观测量 \(f,g\in L^2(\pi^{(q)})\)，若 \((X_n)_{n\ge 0}\) 为以 \(\mathcal{P}^{(q)}\) 为转移核、以 \(\pi^{(q)}\) 为初始分布的平稳链，则有谱控制的相关衰减：
\begin{equation}\label{eq:pom-replica-softcore-cov-decay}
\bigl|\mathrm{Cov}_{\pi^{(q)}}(f(X_0),g(X_n))\bigr|
\;\le\;
\|f\|_{L^2(\pi^{(q)})}\,\|g\|_{L^2(\pi^{(q)})}\,
\Bigl(\frac{\Lambda_q}{\rho_q}\Bigr)^{n},
\end{equation}
其中 \(\Lambda_q:=\max\{|\lambda|:\lambda\in\mathrm{spec}(T^{(q)}),\ |\lambda|<\rho_q\}\)。
若 \(f,g\) 进一步限制为 \(S_q\)-对称可观测量（即仅依赖 Hamming 权），则上式中的 \(\Lambda_q\) 可替换为
\(
\max_{1\le i\le q}|\lambda_i^{(q)}|
\)；
并由定理 \ref{thm:pom-replica-softcore-second-exceptional-eigenvalue-asymptotics}--\ref{thm:pom-replica-softcore-gap-ratio-log-separation} 可知当 \(q\to\infty\) 时该量由 \(\nu_q\) 饱和。
结合定理 \ref{thm:pom-replica-softcore-gap-ratio-log-separation} 与命题 \ref{prop:pom-multiplicity-composition-replica-explicit-eigs} 的固定谱簇上界，
可得当 \(q\to\infty\) 时衰减基数满足
\[
\frac{\Lambda_q}{\rho_q}
=\Big(\frac{\varphi}{2}\Big)^q\bigl(1+o(1)\bigr),
\]
从而相关长度在 \(q\) 上呈 \(1/q\) 量级。
\leanverified{paper\_pom\_replica\_doob\_reversible\_markov\_seeds}
\end{theorem}
\begin{proof}[证明要点]
由 Perron--Frobenius 理论与 Doob 变换的标准构造，\(\mathcal{P}^{(q)}\) 为马尔可夫核且以 \(\pi^{(q)}\) 为平稳分布。
又因 \(T^{(q)}\) 对称，\(\mathcal{P}^{(q)}\) 在 \(L^2(\pi^{(q)})\) 上自伴随，且
\(
\mathcal{P}^{(q)}
\)
与
\(
T^{(q)}/\rho_q
\)
相似，故其非平凡谱半径为 \(\Lambda_q/\rho_q\)。
对中心化可观测量应用谱分解即得 \eqref{eq:pom-replica-softcore-cov-decay}。
最后由命题 \ref{prop:pom-multiplicity-composition-replica-explicit-eigs} 的固定谱簇
\(\bigl\{ \frac{(-1)^j\varphi^{q-2j}}{2}\bigr\}_{j\ge 1}\)
给出 \(\Lambda_q\ge \varphi^{q-2}/2\)，并与定理 \ref{thm:pom-replica-softcore-gap-ratio-log-separation} 的例外块主指数律合并即可得到 \(q\to\infty\) 的指数率。证毕。
\end{proof}

\begin{theorem}[固定谱簇的显式张量本征基：\texorpdfstring{$K^{\otimes q}$}{K^⊗q} 的坐标化对角化]\label{thm:pom-replica-softcore-tensor-eigenbasis}
令
\(
K=\begin{pmatrix}1&1\\[2pt]1&0\end{pmatrix}
\),
并取其特征向量
\(
v^+=(\varphi,1)^{\mathsf T}
\)、\(
v^-=(1,-\varphi)^{\mathsf T}
\)，对应特征值 \(\varphi\) 与 \(-\varphi^{-1}\)。
对任意 \(S\subseteq[q]\)，定义张量向量
\[
f_S:=\bigotimes_{i=1}^{q} w_i,\qquad
w_i=
\begin{cases}
v^-,& i\in S,\\
v^+,& i\notin S.
\end{cases}
\]
则对 \(D^{(q)}:=K^{\otimes q}\) 有严格本征关系
\[
D^{(q)}f_S = (-1)^{|S|}\varphi^{\,q-2|S|}\,f_S.
\]
并且由于
\(
T^{(q)}=\tfrac12(\mathbf{1}\mathbf{1}^{\mathsf T}+D^{(q)})
\)
且 \(\mathbf{1}\mathbf{1}^{\mathsf T}\) 在 \(\mathbf{1}^\perp\) 上恒为零，
任取 \(|S|=j\) 的等特征空间内与 \(\mathbf{1}\) 正交的 \(\binom{q}{j}-1\) 维子空间，
皆为 \(T^{(q)}\) 的显式本征子空间，特征值为
\(
\tfrac12(-1)^j\varphi^{q-2j}
\)。
该结论给出除例外谱外全部高重数谱层的完全坐标化描述。
\end{theorem}
\begin{proof}
前一条由 \(D^{(q)}\) 的张量谱乘法性直接给出。
后一条由
\(
T^{(q)}|_{\mathbf{1}^\perp}=\frac12 D^{(q)}|_{\mathbf{1}^\perp}
\)
以及每个 \(D^{(q)}\) 等特征空间（维 \(\binom{q}{j}\)）与 \(\mathbf{1}^\perp\) 的维数估计得到。证毕。
\end{proof}

\begin{theorem}[特征多项式的结构性分解：例外因子与固定因子完全分离]\label{thm:pom-replica-softcore-characteristic-factorization}
记例外块（维数 \(q+1\)）的特征多项式为
\(
\chi_{\mathrm{exc}}^{(q)}(x):=\prod_{i=0}^{q}(x-\lambda_i^{(q)})
\)。
则 \(T^{(q)}\) 的全特征多项式满足严格因子分解
\begin{equation}\label{eq:pom-replica-softcore-characteristic-factorization}
\chi_{T^{(q)}}(x)
=\chi_{\mathrm{exc}}^{(q)}(x)\;
\prod_{j=0}^{q}\Bigl(x-\frac{(-1)^j\varphi^{q-2j}}{2}\Bigr)^{\binom{q}{j}-1}.
\end{equation}
其中固定因子完全由张量谱簇给出，例外因子容纳全部秩一耦合效应；两者在代数上严格解耦。
\leanverified{paper\_pom\_replica\_characteristic\_factorization\_seeds}
\end{theorem}
\begin{proof}
由命题 \ref{prop:pom-multiplicity-composition-replica-explicit-eigs} 及定理 \ref{thm:pom-replica-softcore-tensor-eigenbasis}，
\(\frac{(-1)^j\varphi^{q-2j}}{2}\)（\(0\le j\le q\)）在 \(\mathbf{1}^\perp\) 上的本征重数至少为 \(\binom{q}{j}-1\)，
从而固定谱簇贡献的特征因子为
\(
\prod_{j=0}^{q}(x-\frac{(-1)^j\varphi^{q-2j}}{2})^{\binom{q}{j}-1}
\)。
其余本征值至多 \(q+1\) 个且均落在 \(S_q\)-对称子空间，并与 \(A^{(q)}\) 的谱一致（命题 \ref{prop:pom-multiplicity-composition-sq-symmetry-reduction}），
故剩余因子恰为 \(\chi_{\mathrm{exc}}^{(q)}\)。
将全部本征值按代数重数组合即得 \eqref{eq:pom-replica-softcore-characteristic-factorization}。证毕。
\end{proof}

\paragraph{二字母词--张量迹：软核链幂迹与谱幂和的统一表达}
沿用 \S\ref{subsec:pom-multiplicity-composition-exact-update-softcore} 的记号：\(U_q=\{0,1\}^q\)，
传递矩阵 \(T^{(q)}\in\RR^{U_q\times U_q}\) 由
\(
T^{(q)}_{u,v}=2^{-\mathbf{1}_{\{u\wedge v\neq 0\}}}
\)
定义，且在自然张量基下满足分解
\(
T^{(q)}=\tfrac12(\mathbf{1}\mathbf{1}^{\mathsf T}+K^{\otimes q})
\)
（命题 \ref{prop:pom-multiplicity-composition-replica-explicit-eigs}），其中
\(
K=\bigl(\begin{smallmatrix}1&1\\[2pt]1&0\end{smallmatrix}\bigr)
\)。
本段给出一种与谱分解正交的幂迹表达：把 \(\Tr((T^{(q)})^m)\) 与例外块幂和同时写成二字母词迹的 \(q\) 次幂平均，并由此导出固定谱簇幂和的“单项塌缩”。

\begin{theorem}[二字母词--张量迹公式与例外块幂和]\label{thm:pom-replica-softcore-word-trace-power-sums}
\leanverified{paper\_pom\_replica\_softcore\_word\_trace\_power\_sums}
令 \(E=\RR^2\)，并记
\[
J_2:=\begin{pmatrix}1&1\\[2pt]1&1\end{pmatrix}
=uu^{\mathsf T},
\qquad
u:=(1,1)^{\mathsf T}.
\]
在同构 \(\RR^{U_q}\cong E^{\otimes q}\) 下有 \(\mathbf{1}\mathbf{1}^{\mathsf T}=J_2^{\otimes q}\) 且 \(K^{\otimes q}=(K)^{\otimes q}\)。
对整数 \(m\ge 1\) 定义二字母词集合 \(\mathcal{W}_m:=\{K,J_2\}^m\)。
对任意 \(w=w_1\cdots w_m\in\mathcal{W}_m\) 记
\[
M_w:=w_1w_2\cdots w_m\in\ZZ^{2\times 2},
\qquad
\tau(w):=\Tr(M_w).
\]
则对任意整数 \(q\ge 1\) 成立：
\begin{enumerate}
  \item \textbf{（全空间幂迹的词分解）}
  \begin{equation}\label{eq:pom-replica-softcore-word-trace-full}
  \boxed{\;
  \Tr\!\bigl((T^{(q)})^{m}\bigr)
  =2^{-m}\sum_{w\in\mathcal{W}_m}\tau(w)^{q}\; } .
  \end{equation}

  \item \textbf{（例外块幂和的 \(\mathrm{Sym}^q\) 读出）}
  记例外块谱为 \(\{\lambda_i^{(q)}\}_{i=0}^{q}\)（命题 \ref{prop:pom-multiplicity-composition-sq-symmetry-reduction}），并定义其幂和
  \[
  S_m^{\mathrm{exc}}(q):=\sum_{i=0}^{q}\bigl(\lambda_i^{(q)}\bigr)^{m}
  =\Tr\!\bigl((A^{(q)})^m\bigr).
  \]
  则有
  \begin{equation}\label{eq:pom-replica-softcore-word-trace-exc}
  \boxed{\;
  S_m^{\mathrm{exc}}(q)
  =2^{-m}\sum_{w\in\mathcal{W}_m}\Tr\!\bigl(\mathrm{Sym}^q(M_w)\bigr)\; } .
  \end{equation}

  \item \textbf{（固定谱簇幂和的单项塌缩）}
  令 \(L_m:=\Tr(K^m)\)（Lucas 数），并定义
  \[
  \Omega_m(q):=\Tr\!\bigl(\mathrm{Sym}^q(K^m)\bigr).
  \]
  则
  \begin{equation}\label{eq:pom-replica-softcore-word-trace-collapse}
  \boxed{\;
  \Tr\!\bigl((T^{(q)})^{m}\bigr)-S_m^{\mathrm{exc}}(q)
  =2^{-m}\bigl(L_m^{q}-\Omega_m(q)\bigr)\; } .
  \end{equation}
  等价地，固定谱簇（即 \(\mathbf{1}^\perp\) 上的全部本征值，参见定理 \ref{thm:pom-replica-softcore-characteristic-factorization}）的 \(m\) 次幂和仅由单个 Lucas 幂 \(L_m^q\) 与单项修正 \(\Omega_m(q)\) 决定。
\end{enumerate}
\end{theorem}
\begin{proof}
由 \(\mathbf{1}\mathbf{1}^{\mathsf T}=J_2^{\otimes q}\) 与 \(T^{(q)}=\tfrac12(\mathbf{1}\mathbf{1}^{\mathsf T}+K^{\otimes q})\)，得
\[
T^{(q)}=\frac12\bigl(J_2^{\otimes q}+K^{\otimes q}\bigr).
\]
利用 Kronecker 幂的乘法性 \((X^{\otimes q})(Y^{\otimes q})=(XY)^{\otimes q}\)，展开得到
\[
(T^{(q)})^m
=2^{-m}\sum_{w\in\mathcal{W}_m}(M_w)^{\otimes q}.
\]
对迹，使用 \(\Tr(X\otimes Y)=\Tr(X)\Tr(Y)\) 归纳推出
\(\Tr((M_w)^{\otimes q})=\Tr(M_w)^q=\tau(w)^q\)，从而得到 \eqref{eq:pom-replica-softcore-word-trace-full}。
\par
对 \eqref{eq:pom-replica-softcore-word-trace-exc}，注意 \(S_q\) 对 \(E^{\otimes q}\) 的坐标置换作用之不变子空间与 \(\mathrm{Sym}^q(E)\) 同构；
并且 \(\mathrm{Sym}^q\) 为表示函子：\((M_w)^{\otimes q}\) 在 \(\mathrm{Sym}^q(E)\) 上的限制同构于 \(\mathrm{Sym}^q(M_w)\)。
对上式在 \(\mathrm{Sym}^q(E)\) 上取迹即得 \eqref{eq:pom-replica-softcore-word-trace-exc}。
\par
最后证明塌缩式 \eqref{eq:pom-replica-softcore-word-trace-collapse}。
由于 \(\det(J_2)=0\) 且 \(\det(K)=-1\)，除唯一纯词 \(w=K^m\) 外，任何 \(w\in\mathcal{W}_m\) 只要含一次 \(J_2\) 即满足 \(\det(M_w)=0\)。
对任意 \(2\times 2\) 矩阵 \(M\) 若 \(\det(M)=0\)，则其特征多项式为 \(x(x-\Tr(M))\)，故 \(\mathrm{Sym}^q(M)\) 的全部特征值除 \(\Tr(M)^q\) 外皆为 \(0\)，从而
\[
\Tr\!\bigl(\mathrm{Sym}^q(M)\bigr)=\Tr(M)^q.
\]
因此在 \eqref{eq:pom-replica-softcore-word-trace-full} 与 \eqref{eq:pom-replica-softcore-word-trace-exc} 的两条词和中，除 \(w=K^m\) 一项外逐项相同；
而该唯一差异项分别给出 \(L_m^q=\Tr(K^m)^q\) 与 \(\Omega_m(q)=\Tr(\mathrm{Sym}^q(K^m))\)。
相减即得 \eqref{eq:pom-replica-softcore-word-trace-collapse}。证毕。
\end{proof}

\begin{proposition}[含 \(J_2\) 之二字母词迹的 Fibonacci 乘积分解]\label{prop:pom-replica-softcore-word-trace-fibonacci-factorization}
\leanverified{paper\_pom\_replica\_softcore\_word\_trace\_fibonacci\_factorization}
令 Fibonacci 数列 \((F_n)_{n\ge 0}\) 满足 \(F_0=0,F_1=1\) 与 \(F_{n+2}=F_{n+1}+F_n\)。
取任意 \(m\ge 1\) 与任意包含至少一个 \(J_2\) 的词 \(w\in\mathcal{W}_m\)。
则存在唯一分解
\[
w=K^{a_0}J_2K^{a_1}J_2\cdots J_2K^{a_r},
\qquad
r\ge 1,\ a_i\in\ZZ_{\ge 0},\ \ \sum_{i=0}^{r}a_i+r=m,
\]
并且其迹满足完全因子化公式
\begin{equation}\label{eq:pom-replica-softcore-word-trace-fib-factor}
\boxed{\;
\tau(w)=F_{a_0+a_r+3}\ \prod_{i=1}^{r-1}F_{a_i+3}\; }.
\end{equation}
\end{proposition}
\begin{proof}
唯一分解由字母表 \(\{K,J_2\}\) 的交替分块确定。
由直接归纳可得
\(
K^n=\bigl(\begin{smallmatrix}F_{n+1}&F_n\\[2pt]F_n&F_{n-1}\end{smallmatrix}\bigr)
\)（\(n\ge 1\)），从而
\[
K^n u=\binom{F_{n+2}}{F_{n+1}},
\qquad
u^{\mathsf T}K^n u=F_{n+3}.
\]
利用 \(J_2=uu^{\mathsf T}\) 得
\[
M_w
=K^{a_0}u\cdot (u^{\mathsf T}K^{a_1}u)\cdots (u^{\mathsf T}K^{a_{r-1}}u)\cdot (K^{a_r}u)^{\mathsf T}
=\Bigl(\prod_{i=1}^{r-1}u^{\mathsf T}K^{a_i}u\Bigr)\,(K^{a_0}u)(K^{a_r}u)^{\mathsf T}.
\]
因此
\[
\tau(w)=\Tr(M_w)
=\Bigl(\prod_{i=1}^{r-1}u^{\mathsf T}K^{a_i}u\Bigr)\,(K^{a_r}u)^{\mathsf T}(K^{a_0}u)
=\Bigl(\prod_{i=1}^{r-1}F_{a_i+3}\Bigr)\,u^{\mathsf T}K^{a_0+a_r}u,
\]
代入 \(u^{\mathsf T}K^{a_0+a_r}u=F_{a_0+a_r+3}\) 即得 \eqref{eq:pom-replica-softcore-word-trace-fib-factor}。证毕。
\end{proof}

\begin{theorem}[整数底数簇的循环生成函数闭式]\label{thm:pom-replica-softcore-word-trace-cycle-gf}
固定整数 \(q\ge 1\)。
对 \(m\ge 1\) 定义
\[
W_m(q):=\sum_{\substack{w\in\mathcal{W}_m\\ w\neq K^m}}\tau(w)^q
=\sum_{\substack{w\in\mathcal{W}_m\\ w\ \text{含至少一 }J_2}}\tau(w)^q.
\]
令
\[
f_n:=F_{n+3}^{\,q}\quad(n\ge 0),
\qquad
G_q(z):=\sum_{n\ge 0}f_n z^n=\sum_{n\ge 0}F_{n+3}^{\,q}z^n.
\]
则 \(W_m(q)\) 的普通生成函数满足闭式
\begin{equation}\label{eq:pom-replica-softcore-word-trace-Wm-gf}
\boxed{\;
\sum_{m\ge 1}W_m(q)\,z^m
=\frac{z\,\dfrac{d}{dz}\bigl(zG_q(z)\bigr)}{1-zG_q(z)}
=-z\,\frac{d}{dz}\log\bigl(1-zG_q(z)\bigr)\;}
\end{equation}
（在形式幂级数意义下成立，且在 \(|z|\) 足够小时亦为解析恒等式）。
\end{theorem}
\begin{proof}
由命题 \ref{prop:pom-replica-softcore-word-trace-fibonacci-factorization}，对含至少一 \(J_2\) 的词，
写成 \(w=K^{a_0}J_2K^{a_1}\cdots J_2K^{a_r}\)（\(r\ge 1\)）后有
\[
\tau(w)^q
=f_{a_0+a_r}\prod_{i=1}^{r-1}f_{a_i},
\qquad
\sum_{i=0}^{r}a_i=m-r.
\]
因此
\[
W_m(q)
=\sum_{r=1}^{m}\ \sum_{\substack{a_0,\dots,a_r\ge 0\\ a_0+\cdots+a_r=m-r}}
f_{a_0+a_r}\prod_{i=1}^{r-1}f_{a_i}.
\]
将其乘以 \(z^m\) 并对 \(m\ge 1\) 求和、交换求和次序得
\[
\sum_{m\ge 1}W_m(q)z^m
=\sum_{r\ge 1}z^r\Bigl(\sum_{a_0,a_r\ge 0} f_{a_0+a_r}z^{a_0+a_r}\Bigr)\,G_q(z)^{r-1}.
\]
注意
\(
\sum_{a_0,a_r\ge 0} f_{a_0+a_r}z^{a_0+a_r}
=\sum_{s\ge 0}(s+1)f_s z^s
=\frac{d}{dz}\bigl(zG_q(z)\bigr)
\)，
于是对 \(r\) 求和得到 \eqref{eq:pom-replica-softcore-word-trace-Wm-gf}。证毕。
\end{proof}
\leanverified{paper\_pom\_replica\_softcore\_word\_trace\_cycle\_gf}

\begin{proposition}[固定 \(K\) 计数下的词迹极值]\label{prop:pom-replica-softcore-word-trace-extremal}
\leanverified{paper\_pom\_replica\_softcore\_word\_trace\_extremal}
令 \(m\ge 2\) 且 \(0\le k\le m-1\)。
在所有长度为 \(m\) 且含至少一个 \(J_2\) 的词中，进一步固定字母 \(K\) 的出现次数为 \(k\)：
\[
\mathcal{W}_{m,k}:=\bigl\{w\in\mathcal{W}_m:\ \#\{i:\ w_i=K\}=k,\ \ w\ \text{含至少一 }J_2\bigr\}.
\]
则迹满足上确界
\begin{equation}\label{eq:pom-replica-softcore-word-trace-extremal}
\boxed{\;
\max_{w\in\mathcal{W}_{m,k}}\tau(w)=2^{m-k-1}F_{k+3}\;}
\end{equation}
并且等号当且仅当词中全部 \(K\) 集中为一段连续块（等价地，在分解 \(w=K^{a_0}J_2K^{a_1}\cdots J_2K^{a_r}\) 中至多一个 \(a_i\) 非零）。
特别地，排除唯一最大词 \(J_2^m\)（其迹为 \(2^m\)）后，有
\[
\max_{\substack{w\in\mathcal{W}_m\\ w\neq J_2^m\\ w\ \text{含至少一 }J_2}}\tau(w)=3\cdot 2^{m-2},
\]
且恰由“仅含一个 \(K\)”的 \(m\) 个词达到。
\end{proposition}
\begin{proof}
先证对任意 \(a,b\ge 0\) 有不等式
\begin{equation}\label{eq:pom-replica-softcore-fib-merge-ineq}
F_{a+3}F_{b+3}\le 2F_{a+b+3}.
\end{equation}
固定 \(a\)，令 \(g_b:=2F_{a+b+3}-F_{a+3}F_{b+3}\)。
由 Fibonacci 递推可得 \(g_{b+2}=g_{b+1}+g_b\)；
并且 \(g_0=0\)、\(g_1=F_{a+1}\ge 0\)，故对所有 \(b\ge 0\) 有 \(g_b\ge 0\)，即得 \eqref{eq:pom-replica-softcore-fib-merge-ineq}。
\par
现取 \(w\in\mathcal{W}_{m,k}\)，写其唯一分解 \(w=K^{a_0}J_2K^{a_1}\cdots J_2K^{a_r}\)，其中 \(r=m-k\ge 1\) 且 \(\sum_i a_i=k\)。
由 \eqref{eq:pom-replica-softcore-word-trace-fib-factor}，
\[
\tau(w)=F_{a_0+a_r+3}\prod_{i=1}^{r-1}F_{a_i+3}.
\]
对内部乘积反复应用 \eqref{eq:pom-replica-softcore-fib-merge-ineq} 将其并入一个因子，得到
\[
\prod_{i=1}^{r-1}F_{a_i+3}\le 2^{r-2}F_{a_1+\cdots+a_{r-1}+3}.
\]
再对 \(F_{a_0+a_r+3}F_{a_1+\cdots+a_{r-1}+3}\) 应用一次 \eqref{eq:pom-replica-softcore-fib-merge-ineq}，得到
\[
\tau(w)\le 2^{r-1}F_{k+3}=2^{m-k-1}F_{k+3},
\]
即 \eqref{eq:pom-replica-softcore-word-trace-extremal}。
等号成立当且仅当上述每次并入都取等号；由 \(g_0=0\) 可知 \eqref{eq:pom-replica-softcore-fib-merge-ineq} 取等号当且仅当其中一个参数为 \(0\)，
从而必须使 \(a_i\) 除至多一个外皆为 \(0\)，即全部 \(K\) 集中为单块。
\par
最后，对 \(w\neq J_2^m\) 的情形，最大值对应 \(k=1\)：
\(
2^{m-2}F_4=3\cdot 2^{m-2}
\)，
并且当且仅当唯一的 \(K\) 位于 \(m\) 个位置之一时取到，故达到词数为 \(m\)。
证毕。
\end{proof}

\begin{corollary}[词迹底数的最小值与精确重数]\label{cor:pom-replica-softcore-word-trace-minimal-base-multiplicity}
\leanverified{paper\_pom\_replica\_softcore\_word\_trace\_minimal\_base\_multiplicity}
固定整数 \(m\ge 1\) 并沿用 \(\mathcal W_m=\{K,J_2\}^m\) 与 \(\tau(w)=\Tr(M_w)\) 的记号。
在全部含至少一个 \(J_2\) 的词迹集合
\[
\bigl\{\tau(w):w\in\mathcal W_m,\ w\neq K^m\bigr\}
\]
中，最小值满足
\[
\boxed{\ \min_{\substack{w\in\mathcal W_m\\ w\neq K^m}}\tau(w)=F_{m+2}\ }
\]
且达到该最小值的词恰有 \(m\) 个，即全部“仅含一个 \(J_2\)”的词
\(
K^{a}J_2K^{m-1-a}
\)
（\(0\le a\le m-1\)）。
\end{corollary}
\begin{proof}
取任意 \(w\in\mathcal W_m\) 且含至少一个 \(J_2\)。
由于 \(\tau(w)=\Tr(M_w)\) 在循环移位下不变，可不失一般性地把 \(w\) 循环移位为以 \(J_2\) 起始的表示：
\[
w=J_2K^{g_1}J_2K^{g_2}\cdots J_2K^{g_r},
\qquad r\ge 1,\ g_j\in\ZZ_{\ge 0},\ \sum_{j=1}^{r}g_j=m-r.
\]
将该表示代入命题 \ref{prop:pom-replica-softcore-word-trace-fibonacci-factorization}（取 \(a_0=0\) 且 \(a_j=g_j\)）得到
\[
\tau(w)=\prod_{j=1}^{r}F_{g_j+3}.
\]
利用 Fibonacci 恒等式 \(F_pF_q=F_{p+q-2}+F_{p-2}F_{q-2}\)（\(p,q\ge 2\)），可得
\(
F_pF_q\ge F_{p+q-2}
\)
且当 \(p,q\ge 3\) 时为严格不等式。
将该不等式迭代合并上述乘积，得到
\[
\tau(w)\ \ge\ F_{\,(g_1+3)+\cdots+(g_r+3)-2(r-1)}=F_{m+2},
\]
并且当 \(r\ge 2\) 时严格大于 \(F_{m+2}\)。
因此最小值只能在 \(r=1\) 时取得，且此时 \(\tau(w)=F_{(m-1)+3}=F_{m+2}\)。
含恰一个 \(J_2\) 的词由 \(J_2\) 的位置唯一确定，共有 \(m\) 个，结论成立。证毕。
\end{proof}

\endinput

\paragraph{温度参数延拓与 Bernoulli 子图矩：一般图配分函数、Hankel 正性与例外块端点刚性}
在软核复制链口径中，传递矩阵 \(T^{(q)}\) 对应于“重叠边”代价 \(1/2\) 的特例。
本段将该代价推广为温度门参数 \(p\in[0,1]\)，并把环幂迹与一般图上的配分函数严格识别为 Bernoulli 随机子图之独立集计数的 \(q\)-次矩。
该识别一方面在 \(q\)-轴上自动给出 Stieltjes 矩结构与 Hankel 正半定，另一方面在 \(m\)-轴上诱导一个带权的二字母词迹展开，并保持“固定谱簇幂和的单项塌缩”机制在全参数族内的稳定性。

\begin{definition}[温度门软核与 Bernoulli 子图矩]\label{def:pom-replica-softcore-temperature-kernel-bernoulli}
取整数 \(q\ge 1\)，令 \(U_q:=\{0,1\}^q\)，并以 \(u\wedge v\) 表示逐坐标与。
对 \(p\in[0,1]\) 定义温度门软核 \(T^{(q)}_p\in\RR^{U_q\times U_q}\)：
\begin{equation}\label{eq:pom-replica-softcore-temperature-kernel}
\boxed{\;
T^{(q)}_p(u,v)
:=
\begin{cases}
1,& u\wedge v=0,\\[2pt]
p,& u\wedge v\neq 0,
\end{cases}
\qquad(u,v\in U_q).\;}
\end{equation}
等价地，令 \(\mathbf{1}\) 为全 \(1\) 向量并令 \(K=\bigl(\begin{smallmatrix}1&1\\[2pt]1&0\end{smallmatrix}\bigr)\)，则有张量分解
\begin{equation}\label{eq:pom-replica-softcore-temperature-decomposition}
\boxed{\;
T^{(q)}_p
=p\,\mathbf{1}\mathbf{1}^{\mathsf T}+(1-p)\,K^{\otimes q}\;}
\end{equation}
（其中 \((K^{\otimes q})_{u,v}=\mathbf{1}_{\{u\wedge v=0\}}\)）。
当 \(p=\tfrac12\) 时，\eqref{eq:pom-replica-softcore-temperature-kernel} 退化为命题 \ref{prop:pom-multiplicity-composition-replica-softcore-transfer} 的软核 \(T^{(q)}\)。
\par
对任意有限简单图 \(G=(V,E)\)，定义 \(q\)-复制配分函数
\begin{equation}\label{eq:pom-replica-softcore-graph-partition-function}
\boxed{\;
Z_G^{(q)}(p)
:=\sum_{\sigma:V\to U_q}\ \prod_{\{x,y\}\in E}T^{(q)}_p\bigl(\sigma(x),\sigma(y)\bigr)\;}
\end{equation}
并定义独立集计数
\[
\mathsf{I}(H)
:=\#\Bigl\{\tau:V(H)\to\{0,1\}:\ \tau(x)\tau(y)=0\ \text{对所有 }\{x,y\}\in E(H)\Bigr\}.
\]
\end{definition}

\begin{remark}[单比特温度核的 NAND 表述]\label{rem:pom-replica-softcore-temperature-kernel-nand}
在定义 \ref{def:pom-replica-softcore-temperature-kernel-bernoulli} 中，单比特矩阵 \(K\) 满足
\[
K_{a,b}=1\ \Longleftrightarrow\ a\wedge b=0
\ \Longleftrightarrow\ (a,b)\neq(1,1),
\qquad a,b\in\{0,1\}.
\]
因此 \(K\) 与布尔 NAND 门（仅在 \((1,1)\) 处输出 \(0\)）的真值表矩阵同构；其张量幂 \(K^{\otimes q}\) 则编码 \(q\) 坐标同时满足 \(u\wedge v=0\) 的约束核。
\end{remark}

\begin{theorem}[指数温度调度下的 Perron 自由能相变律]\label{thm:pom-replica-softcore-temperature-free-energy-phase-transition}
\leanverified{paper\_pom\_replica\_softcore\_temperature\_free\_energy\_phase\_transition}
令 \(\varphi=\frac{1+\sqrt5}{2}\)，并令 \(\rho(\cdot)\) 表示谱半径。
对任意常数 \(\alpha\ge 0\)，取温度参数随阶数指数衰减
\[
p_q:=2^{-\alpha q}\qquad(q\in\NN).
\]
则极限
\[
F(\alpha):=\lim_{q\to\infty}\frac1q\log \rho\!\bigl(T_{p_q}^{(q)}\bigr)
\]
存在，且满足闭式
\[
F(\alpha)=\max\bigl\{\log\varphi,\ (1-\alpha)\log 2\bigr\}.
\]
等价地，存在临界指数
\[
\alpha_c:=1-\log_2\varphi
\]
使得当 \(\alpha<\alpha_c\) 时 \(F(\alpha)=(1-\alpha)\log 2\)，当 \(\alpha\ge\alpha_c\) 时 \(F(\alpha)=\log\varphi\)。
\end{theorem}

\begin{proof}
由 \eqref{eq:pom-replica-softcore-temperature-decomposition}，
\(
T^{(q)}_{p}=p\,\mathbf1\mathbf1^{\mathsf T}+(1-p)K^{\otimes q}
\)。
注意到 \(\mathbf1\mathbf1^{\mathsf T}\) 为秩一非负矩阵，其 Perron 根为 \(\mathbf1^{\mathsf T}\mathbf1=2^q\)，故
\[
\rho\!\bigl(p\,\mathbf1\mathbf1^{\mathsf T}\bigr)=p\,2^q.
\]
又由于张量谱乘法性与 \(\rho(K)=\varphi\)，有
\[
\rho(K^{\otimes q})=\rho(K)^q=\varphi^q,
\qquad
\rho\!\bigl((1-p)K^{\otimes q}\bigr)=(1-p)\varphi^q.
\]
并且按元素非负性有 \(T_p^{(q)}\ge p\,\mathbf1\mathbf1^{\mathsf T}\) 与 \(T_p^{(q)}\ge(1-p)K^{\otimes q}\)，故由谱半径对非负矩阵的单调性得到下界
\[
\rho(T_p^{(q)})\ \ge\ \max\bigl\{p\,2^q,\ (1-p)\varphi^q\bigr\}.
\]
另一方面，\(T_p^{(q)}\) 为对称非负矩阵，故 \(\rho(T_p^{(q)})=\|T_p^{(q)}\|_2\)。
由算子范数三角不等式得
\[
\rho(T_p^{(q)})
\le
\bigl\|p\,\mathbf1\mathbf1^{\mathsf T}\bigr\|_2+\bigl\|(1-p)K^{\otimes q}\bigr\|_2
=p\,2^q+(1-p)\varphi^q.
\]
因此对任意 \(p\in[0,1]\) 有夹逼
\[
\max\bigl\{p\,2^q,\ (1-p)\varphi^q\bigr\}
\ \le\
\rho(T_p^{(q)})
\ \le\
p\,2^q+(1-p)\varphi^q.
\]
取 \(p=p_q=2^{-\alpha q}\) 得 \(p_q\,2^q=2^{(1-\alpha)q}\)，且 \((1-p_q)\varphi^q=\varphi^q(1+o(1))\)。
由两项指数增长的最大指数主导律，
\[
\lim_{q\to\infty}\frac1q\log\bigl(2^{(1-\alpha)q}+\varphi^q\bigr)
=\max\bigl\{(1-\alpha)\log 2,\ \log\varphi\bigr\},
\]
结合夹逼即得所示闭式。
临界指数由 \((1-\alpha)\log2=\log\varphi\) 等价变形得到。证毕。
\end{proof}

\begin{theorem}[对称门核的 Perron--Landauer 极限律]\label{thm:pom-perron-landauer-limit-law}
令 \(A\in\RR_{\ge 0}^{d\times d}\) 为实对称非负矩阵，并记 \(\lambda:=\rho(A)\)。
对每个 \(q\ge 1\) 与任意温度调度 \(p_q\in[0,1]\)，在 \(\RR^{d^q}\cong(\RR^{d})^{\otimes q}\) 上定义
\begin{equation}\label{eq:pom-perron-landauer-tensor-soft-kernel}
\boxed{\;
T^{(q)}_{p_q}(A)
:=
p_q\,\mathbf 1\mathbf 1^{\mathsf T}+(1-p_q)\,A^{\otimes q}\;}
\end{equation}
其中 \(\mathbf 1\in\RR^{d^q}\) 为全 \(1\) 向量。
假设指数率极限
\[
L:=\lim_{q\to\infty}\frac{1}{q}\log p_q\ \in[-\infty,0]
\]
存在（约定 \(\log 0=-\infty\)）。
则极限
\[
F(A;\{p_q\})
:=\lim_{q\to\infty}\frac{1}{q}\log\rho\!\bigl(T^{(q)}_{p_q}(A)\bigr)
\]
存在并满足闭式
\[
F(A;\{p_q\})=\max\bigl\{\log\lambda,\ \log d + L\bigr\}.
\]
特别地，若 \(p_q=d^{-\alpha_q q}\) 且 \(\alpha_q\to\alpha\)，则
\[
F(A;\{p_q\})=\max\bigl\{\log\lambda,\ (1-\alpha)\log d\bigr\}.
\]
\end{theorem}
\begin{proof}
由于 \(A\) 对称，\(\rho(A)\) 为最大特征值，且对任意 \(q\) 有 \(\rho(A^{\otimes q})=\rho(A)^q=\lambda^q\)。
又因 \(\mathbf 1\mathbf 1^{\mathsf T}\) 为秩一非负对称矩阵，其唯一非零特征值为 \(\mathbf 1^{\mathsf T}\mathbf 1=d^q\)，故
\[
\rho\!\bigl(p_q\,\mathbf 1\mathbf 1^{\mathsf T}\bigr)=p_q d^q,
\qquad
\rho\!\bigl((1-p_q)A^{\otimes q}\bigr)=(1-p_q)\lambda^q.
\]

由逐项非负性有
\(
T^{(q)}_{p_q}(A)\ge p_q\,\mathbf 1\mathbf 1^{\mathsf T}
\)
与
\(
T^{(q)}_{p_q}(A)\ge (1-p_q)A^{\otimes q}
\)，
从而谱半径的 Perron 单调性给出下界
\[
\rho\!\bigl(T^{(q)}_{p_q}(A)\bigr)\ \ge\ \max\{p_q d^q,\ (1-p_q)\lambda^q\}.
\]
另一方面，\(T^{(q)}_{p_q}(A)\) 仍为实对称矩阵，故 \(\rho(\cdot)=\|\cdot\|_2\)。
由算子范数的三角不等式，
\[
\rho\!\bigl(T^{(q)}_{p_q}(A)\bigr)
\le
\bigl\|p_q\,\mathbf 1\mathbf 1^{\mathsf T}\bigr\|_2+\bigl\|(1-p_q)A^{\otimes q}\bigr\|_2
=p_q d^q+(1-p_q)\lambda^q.
\]
因此对每个 \(q\) 有夹逼
\[
\max\{p_q d^q,\ (1-p_q)\lambda^q\}
\ \le\
\rho\!\bigl(T^{(q)}_{p_q}(A)\bigr)
\ \le\
p_q d^q+(1-p_q)\lambda^q.
\]

将三项取对数并除以 \(q\)。
若 \(p_q=0\) 在无穷多 \(q\) 上出现，则由 \(L=\lim q^{-1}\log p_q\) 的存在性只能是 \(L=-\infty\) 且 \(p_q=0\) 对充分大 \(q\) 恒成立；此时 \(T^{(q)}_{p_q}(A)=A^{\otimes q}\) 对充分大 \(q\) 成立，极限值为 \(\log\lambda=\max\{\log\lambda,\log d+L\}\)，结论平凡。
下文可因此假设 \(p_q>0\) 对充分大 \(q\) 成立。
注意到 \(\lim_{q\to\infty}\frac1q\log(1-p_q)=0\)（因 \(\log(1-p_q)\) 有界且 \(q\to\infty\)），且
\[
\frac1q\log(p_q d^q)=\log d+\frac1q\log p_q\ \to\ \log d+L.
\]
由两项指数主导律
\(
\lim_{q\to\infty}\frac1q\log(x_q+y_q)=\max\{\lim\frac1q\log x_q,\lim\frac1q\log y_q\}
\)
（当两极限存在且 \(x_q,y_q>0\)）以及上述夹逼，即得极限存在并等于 \(\max\{\log\lambda,\log d+L\}\)。
最后一条取 \(p_q=d^{-\alpha_q q}\) 并令 \(L=-\alpha\log d\) 即得。证毕。
\end{proof}



\begin{theorem}[Bernoulli 子图的独立集矩表示]\label{thm:pom-replica-softcore-bernoulli-subgraph-moment-representation}
在定义 \ref{def:pom-replica-softcore-temperature-kernel-bernoulli} 的设定下，对任意有限简单图 \(G=(V,E)\)、任意整数 \(q\ge 1\) 与任意 \(p\in[0,1]\) 有严格恒等式
\begin{equation}\label{eq:pom-replica-softcore-bernoulli-subgraph-moment-representation}
\boxed{\;
Z_G^{(q)}(p)
=\sum_{F\subseteq E}p^{|E|-|F|}(1-p)^{|F|}\,\mathsf{I}(G_F)^{q}\;}
\end{equation}
其中 \(G_F:=(V,F)\) 为生成子图。
等价地，若 \(\mathbf{F}\subseteq E\) 为按边独立取样、以概率 \((1-p)\) 保留边的 Bernoulli 随机子图，则
\[
Z_G^{(q)}(p)=\mathbb{E}\bigl[\mathsf{I}(G_{\mathbf{F}})^{q}\bigr].
\]
\leanverified{paper\_pom\_replica\_bernoulli\_subgraph\_moment\_seeds}
\leanverified{paper\_pom\_replica\_softcore\_bernoulli\_subgraph\_moment\_representation}
\end{theorem}
\begin{proof}
对每条边 \(e=\{x,y\}\in E\) 定义
\(
D_e(\sigma):=\mathbf{1}_{\{\sigma(x)\wedge\sigma(y)=0\}}
\)。
由 \eqref{eq:pom-replica-softcore-temperature-kernel}，有
\(
T^{(q)}_p\bigl(\sigma(x),\sigma(y)\bigr)=p+(1-p)D_e(\sigma)
\)。
因此
\[
\prod_{e\in E}\bigl(p+(1-p)D_e(\sigma)\bigr)
=\sum_{F\subseteq E}p^{|E|-|F|}(1-p)^{|F|}\prod_{e\in F}D_e(\sigma),
\]
将其代回 \eqref{eq:pom-replica-softcore-graph-partition-function} 并交换求和次序得
\[
Z_G^{(q)}(p)
=\sum_{F\subseteq E}p^{|E|-|F|}(1-p)^{|F|}
\sum_{\sigma:V\to U_q}\prod_{e\in F}D_e(\sigma).
\]
对固定 \(F\)，写 \(\sigma=(\sigma^{(1)},\dots,\sigma^{(q)})\) 为 \(q\) 个坐标的二元赋值，其中 \(\sigma^{(a)}:V\to\{0,1\}\)。
注意到
\(
D_e(\sigma)=\prod_{a=1}^q\mathbf{1}_{\{\sigma^{(a)}(x)\sigma^{(a)}(y)=0\}}
\)，
从而
\[
\prod_{e\in F}D_e(\sigma)
=\prod_{a=1}^{q}\ \prod_{\{x,y\}\in F}\mathbf{1}_{\{\sigma^{(a)}(x)\sigma^{(a)}(y)=0\}}.
\]
因此坐标 \(a\) 间完全解耦，并且每个坐标的允许赋值数恰为 \(\mathsf{I}(G_F)\)。
故内层求和为 \(\mathsf{I}(G_F)^q\)，代回即得 \eqref{eq:pom-replica-softcore-bernoulli-subgraph-moment-representation}。证毕。
\end{proof}

\begin{corollary}[环幂迹的 Bernoulli 子图矩表示]\label{cor:pom-replica-softcore-cycle-trace-bernoulli-moment}
取环图 \(C_m\)（\(m\ge 1\)），则对任意整数 \(q\ge 1\) 与任意 \(p\in[0,1]\) 有
\begin{equation}\label{eq:pom-replica-softcore-cycle-trace-bernoulli-moment}
\boxed{\;
\Tr\!\bigl((T^{(q)}_p)^m\bigr)
=\sum_{F\subseteq E(C_m)}p^{m-|F|}(1-p)^{|F|}\,\mathsf{I}\bigl((C_m)_F\bigr)^{q}\;}
\end{equation}
其中 \((C_m)_F\) 为生成子图。
\leanverified{paper\_pom\_replica\_cycle\_trace\_bernoulli\_seeds}
\leanverified{paper\_pom\_replica\_cycle\_trace\_bernoulli\_package}
\end{corollary}
\begin{proof}
将定义 \ref{def:pom-replica-softcore-temperature-kernel-bernoulli} 的 \(Z_G^{(q)}(p)\) 取 \(G=C_m\)。
注意到对环图有矩阵恒等式 \(Z_{C_m}^{(q)}(p)=\Tr((T^{(q)}_p)^m)\)，再应用定理 \ref{thm:pom-replica-softcore-bernoulli-subgraph-moment-representation} 即得结论。证毕。
\end{proof}

\leanverified{paper\_pom\_replica\_softcore\_bernoulli\_stieltjes\_hankel}
\begin{proposition}[Stieltjes 矩结构与 Hankel 正半定]\label{prop:pom-replica-softcore-bernoulli-stieltjes-hankel}
固定有限简单图 \(G=(V,E)\) 与 \(p\in[0,1]\)，令 \(a_q:=Z_G^{(q)}(p)\)（\(q\ge 0\)，并约定 \(Z_G^{(0)}(p)=\sum_{F\subseteq E}p^{|E|-|F|}(1-p)^{|F|}=1\)）。
则存在有限支撑概率测度 \(\mu_{G,p}\) 使
\begin{equation}\label{eq:pom-replica-softcore-bernoulli-stieltjes}
\boxed{\;
a_q=\int_{[0,\infty)} x^{q}\,\dd\mu_{G,p}(x)\qquad(q\ge 0)\;}
\end{equation}
从而 \((a_q)_{q\ge 0}\) 为有限支撑的 Stieltjes 矩序列。
特别地，对任意 \(n\ge 0\)，Hankel 矩阵 \(\bigl[a_{i+j}\bigr]_{0\le i,j\le n}\) 正半定，因而全部主子式非负。
此外，若 \(\mu_{G,p}\) 的支撑至少含两点，则对所有整数 \(q\ge 1\) 有严格对数凸
\[
a_q^2<a_{q-1}a_{q+1}.
\]
\end{proposition}
\begin{proof}
由 \eqref{eq:pom-replica-softcore-bernoulli-subgraph-moment-representation}，
\(
a_q=\sum_{F\subseteq E}w(F)\,\mathsf{I}(G_F)^q
\)，其中 \(w(F)=p^{|E|-|F|}(1-p)^{|F|}\) 且 \(\sum_F w(F)=1\)。
令
\(
\mu_{G,p}:=\sum_{F\subseteq E}w(F)\,\delta_{\mathsf{I}(G_F)}
\)
即得 \eqref{eq:pom-replica-softcore-bernoulli-stieltjes}。
\par
Hankel 正半定性为 Gram 矩阵表述：对任意 \(c_0,\dots,c_n\in\RR\)，有
\[
\sum_{i,j=0}^{n}c_ic_j a_{i+j}
=\int\Bigl(\sum_{i=0}^{n}c_i x^i\Bigr)^2\dd\mu_{G,p}(x)\ge 0.
\]
严格对数凸由 Cauchy--Schwarz 在非退化分布下取严格不等式得到。证毕。
\end{proof}

\begin{proposition}[\texorpdfstring{$S_q$}{Sq}-对称降维的温度延拓]\label{prop:pom-replica-softcore-temperature-sq-reduction}
\leanverified{paper\_pom\_replica\_softcore\_temperature\_sq\_reduction}
固定整数 \(q\ge 1\) 与 \(p\in[0,1]\)。
矩阵 \(T_p^{(q)}\) 在坐标置换群 \(S_q\) 下不变，因此其 Perron 根位于 \(S_q\)-对称子空间。
在以 Hamming 权 \(k=|u|\in\{0,\dots,q\}\) 分层的基下，该限制矩阵为 \(A_p^{(q)}\in\RR^{(q+1)\times(q+1)}\)，其条目为
\begin{equation}\label{eq:pom-replica-softcore-temperature-Apq-def}
\boxed{\;
\bigl(A_p^{(q)}\bigr)_{k,l}
:=p\binom{q}{l}+(1-p)\binom{q-k}{l},
\qquad 0\le k,l\le q.\;}
\end{equation}
并且有严格等式
\begin{equation}\label{eq:pom-replica-softcore-temperature-sq-reduction}
\boxed{\;
\rho\!\bigl(T_p^{(q)}\bigr)=\rho\!\bigl(A_p^{(q)}\bigr).\;}
\end{equation}
\end{proposition}
\begin{proof}
对 \(u\in U_q\) 取 \(k=|u|\)，并取任意仅依赖权的正函数 \(f(u)=g(|u|)\)。
把 \((T_p^{(q)}f)(u)=\sum_{v}T_p^{(q)}(u,v)f(v)\) 按 \(l=|v|\) 分组：
当 \(|v|=l\) 时满足 \(u\wedge v=0\) 的 \(v\) 个数为 \(\binom{q-k}{l}\)，总数为 \(\binom{q}{l}\)。
故
\[
\sum_{\substack{v\in U_q\\ |v|=l}}T_p^{(q)}(u,v)
=\binom{q-k}{l}\cdot 1+\bigl(\binom{q}{l}-\binom{q-k}{l}\bigr)\cdot p
=p\binom{q}{l}+(1-p)\binom{q-k}{l},
\]
即为 \eqref{eq:pom-replica-softcore-temperature-Apq-def}。
由于 \(T_p^{(q)}\) 非负且 \(S_q\)-不变，Perron 根可由对称子空间内的正本征向量实现，故 \eqref{eq:pom-replica-softcore-temperature-sq-reduction} 成立。证毕。
\end{proof}

\begin{theorem}[例外块行列式的 \texorpdfstring{$(1-p)^q$}{(1-p)^q} 刚性]\label{thm:pom-replica-softcore-temperature-exceptional-determinant}
在命题 \ref{prop:pom-replica-softcore-temperature-sq-reduction} 的设定下，对任意整数 \(q\ge 1\) 与任意 \(p\in[0,1]\) 有
\begin{equation}\label{eq:pom-replica-softcore-temperature-exceptional-determinant}
\boxed{\;
\det\!\bigl(A_p^{(q)}\bigr)=(-1)^{\frac{q(q+1)}{2}}\,(1-p)^{q}.\;}
\end{equation}
等价地，若记 \(\mathrm{spec}(A_p^{(q)})=\{\lambda_i^{(q)}(p)\}_{i=0}^{q}\)，则
\[
\prod_{i=0}^{q}\lambda_i^{(q)}(p)=(-1)^{\frac{q(q+1)}{2}}\,(1-p)^{q}.
\]
\leanverified{paper\_pom\_replica\_exceptional\_det\_seeds}
\leanverified{paper\_pom\_replica\_softcore\_temperature\_exceptional\_determinant}
\end{theorem}
\begin{proof}
令 \(\mathbf{1}\in\RR^{q+1}\) 为全 \(1\) 向量，并记 \(b:=(\binom{q}{0},\dots,\binom{q}{q})^{\mathsf T}\)。
再令 \(C^{(q)}\in\RR^{(q+1)\times(q+1)}\) 满足 \(C^{(q)}_{k,l}=\binom{q-k}{l}\)。
由 \eqref{eq:pom-replica-softcore-temperature-Apq-def} 有分解
\[
A_p^{(q)}=p\,\mathbf{1}b^{\mathsf T}+(1-p)\,C^{(q)}.
\]
注意 \(C^{(q)}e_0=\mathbf{1}\)（因 \(C^{(q)}_{k,0}=1\)），故 \((C^{(q)})^{-1}\mathbf{1}=e_0\)；
并且 \(b^{\mathsf T}e_0=\binom{q}{0}=1\)。
对 \((1-p)C^{(q)}\) 应用秩一行列式引理得到
\[
\det(A_p^{(q)})
=(1-p)^{q+1}\det(C^{(q)})\Bigl(1+\frac{p}{1-p}b^{\mathsf T}(C^{(q)})^{-1}\mathbf{1}\Bigr)
=(1-p)^{q}\det(C^{(q)}).
\]
而 \(\det(C^{(q)})=(-1)^{\frac{q(q+1)}{2}}\) 已由 \eqref{eq:pom-replica-softcore-detC} 给出，代回即得 \eqref{eq:pom-replica-softcore-temperature-exceptional-determinant}。证毕。
\end{proof}

\begin{theorem}[温度延拓下的词迹公式与固定谱簇单项塌缩]\label{thm:pom-replica-softcore-temperature-word-trace-collapse}
\leanverified{paper\_pom\_replica\_softcore\_temperature\_word\_trace\_collapse}
沿用定理 \ref{thm:pom-replica-softcore-word-trace-power-sums} 的二字母词记号：\(\mathcal{W}_m=\{K,J_2\}^m\)、\(M_w=w_1\cdots w_m\)、\(\tau(w)=\Tr(M_w)\)。
对 \(w\in\mathcal{W}_m\) 记
\(
N_J(w):=\#\{i:\ w_i=J_2\}
\) 与 \(N_K(w):=\#\{i:\ w_i=K\}=m-N_J(w)\)。
则对任意整数 \(q\ge 1,m\ge 1\) 与任意 \(p\in[0,1]\) 有
\begin{equation}\label{eq:pom-replica-softcore-temperature-word-trace-full}
\boxed{\;
\Tr\!\bigl((T_p^{(q)})^{m}\bigr)
=\sum_{w\in\mathcal{W}_m}p^{N_J(w)}(1-p)^{N_K(w)}\,\tau(w)^{q}\;}
\end{equation}
并且若记例外块幂和为
\(
S_{m,\mathrm{exc}}(q;p):=\Tr\!\bigl((A_p^{(q)})^{m}\bigr)=\sum_{i=0}^{q}\lambda_i^{(q)}(p)^{m}
\)，则成立单项塌缩恒等式
\begin{equation}\label{eq:pom-replica-softcore-temperature-word-trace-collapse}
\boxed{\;
\Tr\!\bigl((T_p^{(q)})^{m}\bigr)-S_{m,\mathrm{exc}}(q;p)
=(1-p)^{m}\Bigl(L_m^{q}-\Omega_m(q)\Bigr)\;}
\end{equation}
其中 \(L_m=\Tr(K^m)\) 为 Lucas 数，而 \(\Omega_m(q)=\Tr(\mathrm{Sym}^q(K^m))\) 如 \eqref{eq:pom-replica-softcore-word-trace-collapse} 所定义。
\end{theorem}
\begin{proof}
由 \eqref{eq:pom-replica-softcore-temperature-decomposition} 与 Kronecker 幂的乘法性展开得
\[
(T_p^{(q)})^m
=\sum_{w\in\mathcal{W}_m}p^{N_J(w)}(1-p)^{N_K(w)}\,(M_w)^{\otimes q},
\]
再取迹并用 \(\Tr((M_w)^{\otimes q})=\Tr(M_w)^q=\tau(w)^q\) 即得 \eqref{eq:pom-replica-softcore-temperature-word-trace-full}。
\par
另一方面，\((M_w)^{\otimes q}\) 在 \(S_q\)-对称子空间的限制同构于 \(\mathrm{Sym}^q(M_w)\)，从而
\[
S_{m,\mathrm{exc}}(q;p)
=\sum_{w\in\mathcal{W}_m}p^{N_J(w)}(1-p)^{N_K(w)}\,\Tr\!\bigl(\mathrm{Sym}^q(M_w)\bigr).
\]
若 \(w\neq K^m\)，则 \(w\) 含至少一个 \(J_2\)，从而 \(\det(M_w)=0\)，并由定理 \ref{thm:pom-replica-softcore-word-trace-power-sums} 的论证得
\(
\Tr(\mathrm{Sym}^q(M_w))=\Tr(M_w)^q=\tau(w)^q
\)。
故两条加权词和逐项相消，仅余 \(w=K^m\) 一项，其权重为 \((1-p)^m\)，并分别贡献 \(L_m^q\) 与 \(\Omega_m(q)\)。
于是得到 \eqref{eq:pom-replica-softcore-temperature-word-trace-collapse}。证毕。
\end{proof}

\endinput


\begin{theorem}[例外特征多项式的 Vieta 刚性：首一降阶项与常数项全闭式]\label{thm:pom-replica-softcore-exceptional-vieta-endpoints}
对 \(q\ge 2\)，例外特征多项式
\(
\chi_{\mathrm{exc}}^{(q)}(x)=\prod_{i=0}^{q}(x-\lambda_i^{(q)})
\)
满足
\begin{equation}\label{eq:pom-replica-softcore-exceptional-vieta-xq}
\chi_{\mathrm{exc}}^{(q)}(x)
=x^{q+1}
-\Bigl(2^{q-1}+\frac{\varphi^{q+2}+(-1)^q\varphi^{-q}}{2(\varphi+2)}\Bigr)x^{q}
+\cdots,
\end{equation}
并且常数项为
\begin{equation}\label{eq:pom-replica-softcore-exceptional-vieta-const}
\chi_{\mathrm{exc}}^{(q)}(0)
=(-1)^{q+1}\prod_{i=0}^{q}\lambda_i^{(q)}
=(-1)^{q+1+\frac{q(q+1)}{2}}\,2^{-q}.
\end{equation}
从而例外因子的首一降阶系数与常数项在闭式意义下完全刚性。
\leanverified{paper\_pom\_replica\_exceptional\_vieta\_endpoints\_seeds}
\end{theorem}
\begin{proof}
由 Vieta 公式，
\(
[x^{q}]\chi_{\mathrm{exc}}^{(q)}(x)=-\sum_{i=0}^{q}\lambda_i^{(q)}
\)
与
\(
\chi_{\mathrm{exc}}^{(q)}(0)=(-1)^{q+1}\prod_{i=0}^{q}\lambda_i^{(q)}
\)。
分别代入定理 \ref{thm:pom-replica-softcore-exceptional-spectrum-trace} 与
定理 \ref{thm:pom-replica-softcore-exceptional-spectrum-product}
即得 \eqref{eq:pom-replica-softcore-exceptional-vieta-xq}--\eqref{eq:pom-replica-softcore-exceptional-vieta-const}。证毕。
\end{proof}

\begin{corollary}[例外特征多项式一次项系数的显式闭式]\label{cor:pom-replica-softcore-exceptional-charpoly-linear-coefficient}
对任意整数 \(q\ge 2\)，例外特征多项式 \(\chi_{\mathrm{exc}}^{(q)}(x)\) 的一次项系数满足
\[
\boxed{\;
[x]\chi_{\mathrm{exc}}^{(q)}(x)
=(-1)^{\frac{q(q+1)}{2}}\,2^{-q}\,\Bigl(2F_{q+1}-(-1)^q\Bigr)\; } ,
\]
其中 \(F_n\) 为 Fibonacci 数列（\(F_0=0,F_1=1\)）。
\end{corollary}
\begin{proof}
由 Vieta 关系，\([x]\chi_{\mathrm{exc}}^{(q)}(x)=(-1)^q e_q\)，其中
\[
e_q=\sum_{i=0}^{q}\ \prod_{\substack{0\le j\le q\\ j\ne i}}\lambda_j^{(q)}
=\Bigl(\prod_{j=0}^{q}\lambda_j^{(q)}\Bigr)\Bigl(\sum_{i=0}^{q}\frac{1}{\lambda_i^{(q)}}\Bigr).
\]
将
\(\prod_{j=0}^{q}\lambda_j^{(q)}=\det(A^{(q)})=(-1)^{\frac{q(q+1)}{2}}2^{-q}\)
（定理 \ref{thm:pom-replica-softcore-exceptional-spectrum-product}）
与
\(\sum_{i=0}^{q}\frac{1}{\lambda_i^{(q)}}=-1+2(-1)^qF_{q+1}\)
（推论 \ref{cor:xi-terminal-replica-softcore-exceptional-reciprocal-sum-fibonacci}）
代入即得所示闭式。证毕。
\end{proof}

\begin{proposition}[例外特征多项式第二 Vieta 系数的闭式]\label{prop:pom-replica-softcore-exceptional-vieta-e2}
对 \(q\ge 2\)，将
\[
\chi_{\mathrm{exc}}^{(q)}(x)
=\prod_{i=0}^{q}(x-\lambda_i^{(q)})
=x^{q+1}-e_1(q)x^{q}+e_2(q)x^{q-1}-\cdots
\]
写为 Vieta 形式。
则第二 Vieta 系数满足严格闭式
\[
\boxed{\;
e_2(q)=\frac14\Bigl(2^{q}F_{q+1}-F_qF_{q+1}-3^{q}\Bigr)\; } .
\]
\leanverified{paper\_pom\_replica\_exceptional\_vieta\_e2\_seeds}
\end{proposition}
\begin{proof}
Newton 恒等式给出
\[
e_2(q)=\frac12\left(\Bigl(\sum_{i=0}^{q}\lambda_i^{(q)}\Bigr)^2-\sum_{i=0}^{q}\bigl(\lambda_i^{(q)}\bigr)^2\right).
\]
由定理 \ref{thm:pom-replica-softcore-exceptional-spectrum-trace} 的证明，有
\(
\sum_i\lambda_i^{(q)}=2^{q-1}+\frac12F_{q+1}
\)；
而结论 \ref{con:xi-terminal-replica-softcore-exceptional-spectrum-square-sum} 给出
\(
\sum_i(\lambda_i^{(q)})^2=4^{q-1}+\frac12\,3^q+\frac14F_{2q+2}
\)。
代入并使用恒等式 \(F_{2q+2}=F_{q+1}\bigl(F_{q+1}+2F_q\bigr)\) 化简，即得所示表达。证毕。
\end{proof}

\begin{conjecture}[例外特征多项式的不可约性与满对称伽罗瓦群]\label{conj:pom-replica-softcore-exceptional-charpoly-irreducible-symmetric-galois}
对任意 \(q\ge 2\)，例外特征多项式 \(\chi_{\mathrm{exc}}^{(q)}(x)\in\QQ[x]\) 在 \(\QQ\) 上不可约，并且
\[
\mathrm{Gal}\!\left(\chi_{\mathrm{exc}}^{(q)}/\QQ\right)\cong S_{q+1}.
\]
\end{conjecture}

上述猜想在有限范围内的可复现实验证书见表 \ref{tab:replica_softcore_exceptional_charpoly_galois_audit}。
% Auto-generated by scripts/exp_replica_softcore_exceptional_charpoly_galois_audit.py
\begin{table}[H]
\centering
\small
\setlength{\tabcolsep}{8pt}
\renewcommand{\arraystretch}{1.12}
\begin{tabular}{rrrr}
\toprule
$q$ & $\deg\chi_{\mathrm{exc}}^{(q)}$ & irreducible over $\QQ$ & $\mathrm{Gal}$ (SymPy, $q\le 5$) \\
\midrule
2 & 3 & yes & $S_3$ \\
3 & 4 & yes & $S_4$ \\
4 & 5 & yes & $S_5$ \\
5 & 6 & yes & $S_6$ \\
6 & 7 & yes & -- \\
7 & 8 & yes & -- \\
8 & 9 & yes & -- \\
9 & 10 & yes & -- \\
10 & 11 & yes & -- \\
\bottomrule
\end{tabular}
\caption{Computational audit for the exceptional characteristic polynomial $\chi_{\mathrm{exc}}^{(q)}$: irreducibility for $2\le q\le 10$ and Galois group for $2\le q\le 5$ (degree $\le 6$, within SymPy support).}
\label{tab:replica_softcore_exceptional_charpoly_galois_audit}
\end{table}


\begin{theorem}[次大例外本征向量的谱层局域：向 \texorpdfstring{$j=0$}{j=0} 层的指数集中]\label{thm:pom-replica-softcore-second-exceptional-eigenvector-localization}
令 \(\nu_q\) 与定理 \ref{thm:pom-replica-softcore-second-exceptional-eigenvalue-asymptotics} 相同，并取其单位本征向量为 \(\phi^{(q)}\)（取于例外块所在的 \((S_q)\)-对称子空间）。
令 \(P_j\) 表示 \(D^{(q)}=K^{\otimes q}\) 的谱层正交投影（对应本征值 \((-1)^j\varphi^{q-2j}\)）。
则存在常数 \(c>0\) 使得
\begin{equation}\label{eq:pom-replica-softcore-second-eigvec-localization}
1-\bigl\|\mathrm{Proj}_{\mathrm{span}\{P_0\mathbf{1}\}}\phi^{(q)}\bigr\|_2^2
\;=\;
O\!\bigl(r^{q}\bigr),
\qquad q\to\infty,
\end{equation}
其中 \(r=\frac{5+2\sqrt5}{10}\) 同 \eqref{eq:pom-replica-softcore-nuq-asymptotics}。
换言之，\(\nu_q\) 的本征态在例外块内以速率 \(r^{q}\) 指数集中到 \(D^{(q)}\) 的最大正谱层方向（对应 \(j=0\)），并与定理 \ref{thm:pom-replica-softcore-second-exceptional-eigenvalue-asymptotics} 的贴近尺度相一致。
\end{theorem}
\begin{proof}[证明要点]
设 \(E_0:=\mathrm{Ran}(P_0)=\mathrm{span}\{(v^+)^{\otimes q}\}\) 为 \(D^{(q)}\) 的最大正谱层。
由定理 \ref{thm:pom-replica-softcore-second-exceptional-eigenvalue-asymptotics}，\(\nu_q\) 与 \(\varphi^q/2\) 的谱距为 \(O(\varphi^q r^{q})\)。
另一方面，除 \(E_0\) 外的全部谱层与 \(\varphi^q\) 至少相差一个常数因子（例如下一层为 \(\varphi^{q-2}\)），
从而由谱投影的标准摄动估计可将 \(\phi^{(q)}\) 在 \(E_0^\perp\) 的分量控制为 \(O(r^{q/2})\)，
即得到 \eqref{eq:pom-replica-softcore-second-eigvec-localization}。证毕。
\end{proof}

\endinput


\begin{remark}[有限窗张量化与一般 \texorpdfstring{$q\ge 0$}{q>=0} 的指数解耦]\label{rem:pom-multiplicity-composition-window-tensorization}
定理 \ref{thm:pom-multiplicity-composition-exact-conditional-iid} 给出对任意实 \(q\ge 0\) 的有限体积精确条件 i.i.d.\ 表示。
由 \(G_q(t)=1/(1-W_q(t))\) 的正系数解析结构（\(W_q\) 在 \((0,1)\) 上严格递增且 \(W_q(\rho(q))=1\) 为简单根，定理 \ref{thm:pom-multiplicity-real-q-pressure}），
可得对每个固定窗口长度 \(m\) 存在常数 \(C_{m,q}<\infty\) 与 \(\eta(q)\in(0,1)\) 使
\[
\Bigl\|\ \mathcal{L}_{P_L^{(q)}}(\ell_1,\dots,\ell_m)-p_q^{\otimes m}\ \Bigr\|_{\mathrm{TV}}
\le C_{m,q}\,\eta(q)^{L}
\qquad(L\to\infty),
\]
其中 \(\eta(q)\) 可取为相应 renewal/转移算子的谱半径比（支配奇点与次支配奇点的模比）。
该指数解耦把“边界抽样”对有限窗口的影响量化为指数小量，并与 \(q=1\) 的显式谱隙常数 \(2/\lambda(1)^2\) 相容（定理 \ref{thm:pom-multiplicity-composition-exponential-mixing} 与命题 \ref{prop:pom-multiplicity-composition-first-part-exact-and-rate}）。
\end{remark}

\endinput
